\documentclass[12pt,a4paper]{ctexart}
\usepackage{amsmath,amssymb,amsfonts,amsthm,mathtools,stmaryrd}
\usepackage{graphicx}
\usepackage{hyperref}
\usepackage{geometry}
\usepackage{booktabs}
\usepackage{float}
\usepackage{tabularx}
\usepackage{tabulary}
\usepackage{tikz}
\usepackage{forest}
\usepackage{listings}
\usepackage{xcolor}
\usepackage{fontspec}
\xeCJKsetup{AutoFallBack = true}
\setCJKfallbackfamilyfont\CJKrmdefault{Malgun Gothic}[Scale=MatchLowercase]
\setCJKfallbackfamilyfont\CJKsfdefault{Malgun Gothic}[Scale=MatchLowercase]
\setCJKfallbackfamilyfont\CJKttdefault{Malgun Gothic}[Scale=MatchLowercase]
\renewcommand{\textsc}[1]{{\upshape\scshape #1}}

\newfontfamily\ipafont{Segoe UI}[Scale=MatchLowercase]
\newcommand{\ipa}[1]{\text{\ipafont #1}}

\newfontfamily\cyrfont{Segoe UI}[Scale=MatchLowercase]
\newcommand{\cyr}[1]{\text{\cyrfont #1}}

\newtheorem{definition}{定义}[section]
\newtheorem{example}{例子}[section]
\newtheorem{theorem}{定理}[section]

\geometry{left=2.5cm,right=2.5cm,top=3cm,bottom=3cm}

\title{\huge 语言学形式化教程\\ \large From Natural Language to Formal Systems}
\author{Language Tutorial Project}
\date{\today}

\begin{document}

\clearpage
\vspace*{\fill}

\begin{center}
{\huge 语言学形式化教程}\\[0.8em]
{\large From Natural Language to Formal Systems}\\[2.5em]
{\large MS Roslyn}\\[0.6em]
{\large \today}
\end{center}

\vspace{2.5em}

\begin{abstract}
\centering
本教程系统介绍语言学的核心领域及其形式化方法，涵盖语言起源与分类、形式语言理论、形态学、句法学、语义学、音系学、计算语言学，并在最后一章集中讨论不同语言之间的差异。通过数学公式、自动机理论和逻辑演算，展示语言学如何像数学一样进行精确的形式化分析，并说明语言差异如何在共享的深层结构上以参数化的方式呈现。
\end{abstract}

\vspace*{\fill}
\clearpage

\tableofcontents
\clearpage

\section{语言学的起源与历史}

\subsection{语言学的定义与研究对象}

语言学(Linguistics)是系统研究人类语言的学科，其研究对象涵盖语言的多个层面：
\begin{itemize}
    \item \textbf{语音(Phonetics)}: 语音的物理与生理属性
    \item \textbf{音系(Phonology)}: 语音在特定语言系统中的组织方式
    \item \textbf{词法(Morphology)}: 词的内部结构与构词规则
    \item \textbf{句法(Syntax)}: 词组合成短语和句子的规则
    \item \textbf{语义(Semantics)}: 语言单位的意义
    \item \textbf{语用(Pragmatics)}: 语言在语境中的使用
\end{itemize}

语言学不仅是人文学科，更是一门交叉科学。它与数学共享形式化方法(集合论、自动机理论、逻辑演算)，与计算机科学共同构成自然语言处理(NLP)的理论基础，与心理学、神经科学共同研究语言的认知与神经基础。正是这种形式化的追求，使语言学区别于一般的人文研究：它试图用精确的符号系统描述语言规律，正如数学描述自然规律。

\subsection{古代语言研究}

\subsubsection{印度传统:波你尼与梵语语法}

公元前4世纪，印度语法学家波你尼(Pāṇini)在《八章书》(Aṣṭādhyāyī)中建立了梵语语法的3,959条规则。这是人类历史上最早的、堪称完整的形式化语法体系。其特点包括：
\begin{itemize}
    \item 采用类似\emph{生成}的方式：从有限规则推导出一切合法形式
    \item 规则间有明确的优先级和相互作用机制
    \item 定义了语素、词干、词缀的操作系统(类似现代的形态音位学)
\end{itemize}

波你尼的规则系统可以用现代形式语言理论的方式表述：
\begin{equation}
    \text{规则: } X \rightarrow Y \quad \text{在条件 } C \text{ 下}
\end{equation}
其规则网络与乔姆斯基层级中上下文敏感的文法规则(见形式语言一章)在思想上有深刻的一致性，因此波你尼被一些学者称为"生成语法的鼻祖"。

\subsubsection{希腊传统:词类与逻辑}

古希腊学者最早将语言与逻辑联系起来。亚里士多德区分了名词(ónoma)和动词(rhēma)，并注意到"真与假"只有通过命题(句子)才能表达，从而将语言研究与逻辑学结合。斯多葛学派进一步区分了"被说出的声音""被表达的内容"与"外部事物"，初步萌芽了后来\emph{能指/所指}的区分。

\subsubsection{中国训诂学:小学传统}

中国古代的"小学"包含文字学、音韵学、训诂学三个分支：
\begin{itemize}
    \item \textbf{文字学}: 分析汉字的构造。许慎《说文解字》(约公元100年)创立了540部首体系，并总结出象形、指事、会意、形声、转注、假借"六书"理论，是第一部分析汉字结构的系统著作。
    \item \textbf{音韵学}: 研究汉语的语音系统，以"反切"注音法、《切韵》《广韵》等韵书为代表。
    \item \textbf{训诂学}: 解释古语词义，《尔雅》是其鼻祖。
\end{itemize}

\subsubsection{阿拉伯传统:西伯维与早期音系学}

公元8世纪，波斯学者西伯维(Sībawayh)在《书》(Al-Kitāb)中对阿拉伯语进行了详尽的描述性分析，尤其是对辅音的音位对立、发音部位和语音环境做了系统的分类。他的工作被公认为早期音系学的典范，比欧洲的类似研究早了一千多年。

\subsection{现代语言学的诞生}

\subsubsection{历史比较语言学(19世纪)}

19世纪的历史比较语言学标志着语言学成为一门独立的经验科学。其核心方法是\emph{比较法}(Comparative Method)：通过比较亲属语言中规律对应的同源词(cognates)，重建它们的祖语。

格里姆定律(Grimm's Law)描述了从原始印欧语(PIE)到日耳曼语族的系统性辅音转移：
\begin{align}
    \text{PIE } *p &\rightarrow \text{Gmc } f \\
    \text{PIE } *t &\rightarrow \text{Gmc } \theta \\
    \text{PIE } *k &\rightarrow \text{Gmc } h
\end{align}

维尔纳定律(Verner's Law)后来解释了格里姆定律的例外(清音变浊擦音取决于原始重音位置)，体现了"规律无例外，例外有规律"的科学原则。由此，历史语言学确立了音变的规律性假设：
\begin{equation}
    \text{如果语言 } L_1, L_2 \text{ 同源，则存在系统的语音对应: } f: \Sigma_{L_1} \rightarrow \Sigma_{L_2}
\end{equation}

\subsubsection{结构主义语言学(20世纪初)}

瑞士语言学家费迪南·德·索绪尔(Ferdinand de Saussure)在《普通语言学教程》中奠定了结构主义的基础：
\begin{itemize}
    \item \textbf{能指(Signifier)}与\textbf{所指(Signified)}: 语音形象与概念之间是任意、约定俗成的关系
    \item \textbf{语言(Langue)}与\textbf{言语(Parole)}: 语言系统 vs. 具体的话语行为
    \item \textbf{共时(Synchronic)}与\textbf{历时(Diachronic)}研究: 系统状态的描写 vs. 历史演变
\end{itemize}

索绪尔的核心洞见是：语言是一个\emph{系统}，其中每个单位的价值由它在系统中的对立关系决定。这一思想后来被形式化为"结构即关系之网络"。

结构主义在布拉格学派(特鲁别茨柯伊、雅柯布森)手中发展为功能结构主义。特鲁别茨柯伊(Trubetzkoy)提出了\emph{音位}(phoneme)概念：一个语言系统中能区分意义的最小语音单位。音位正是通过\emph{对立}(opposition)来定义的：
\begin{equation}
    \text{音位} = \frac{\text{语音特征}}{\text{区别意义的最小对立}}
\end{equation}

美国结构主义(布龙菲尔德 Bloomfield)则强调描写主义与分布分析，主张基于语料对语言进行严格的、可操作的分析。

\subsection{生成语法的革命}

诺姆·乔姆斯基(Noam Chomsky)在1957年《句法结构》中发起了生成语法革命，其核心理念：
\begin{itemize}
    \item 语言具有\textbf{普遍语法}(Universal Grammar, UG)：人类天生的语言官能
    \item 句法是\textbf{生成性}的：有限的规则系统可以生成无限的合法句子
    \item 区分\textbf{能力}(Competence)与\textbf{表现}(Performance)：内在的语言知识 vs. 实际的使用行为
    \item \textbf{刺激贫乏论}(Poverty of the Stimulus)：儿童从有限的输入中习得复杂的语法，说明语言知识有先天基础
\end{itemize}

生成语法把语言学研究从"收集语料、分类描写"转向"建立可以生成一切合法句子、排除一切非法句子的规则系统"。这一目标与形式语言理论完全一致：

\begin{equation}
    G = \langle V_N, V_T, P, S \rangle
\end{equation}
其中 $V_N$ 是非终结符集合，$V_T$ 是终结符集合，$P$ 是产生式规则集合，$S$ 是起始符号。生成语法所生成的语言就是
\begin{equation}
    L(G) = \{ w \in V_T^* \mid S \Rightarrow^* w \}
\end{equation}

\subsubsection{从规则到原则与参数}

生成语法自身也在不断演进：从早期规则系统(短语结构规则+转换规则)，到原则与参数理论(Principles and Parameters)，再到最简方案(Minimal Program)。原则与参数理论认为：人类语言共享一组普遍原则(原则)，而语言之间的差异只是少数\emph{参数}(parameters)取值不同。这为"语言之间的差异"提供了形式化的解释框架，我们将在最后一章详细展开。

\subsection{认知语言学与功能学派}

20世纪后期，学界出现了对"纯形式句法"的反拨：

\subsubsection{认知语言学}

认知语言学强调语言与一般认知能力密不可分：
\begin{itemize}
    \item 语言是\textbf{认知能力}的一部分，而非独立模块
    \item \textbf{隐喻}是思维的基本机制(Lakoff \& Johnson, 1980)：概念隐喻(如"争论是战争")塑造语言表达
    \item 语法是\textbf{用法为基础}(Usage-based)的：语法结构是从具体语言使用中抽象出来的
    \item \textbf{构式语法}(Construction Grammar)：把"词素-词-短语-句式"都看作"形式-意义"配对体
\end{itemize}

\subsubsection{功能语言学}

功能学派(韩礼德 Halliday 的系统功能语法等)强调语言的功能：
\begin{itemize}
    \item 概念功能：语言如何表征世界
    \item 人际功能：语言如何建立社会关系
    \item 语篇功能：语言如何组织信息
\end{itemize}

\subsection{计算语言学时代}

21世纪以来，语言学与计算机科学深度融合：
\begin{itemize}
    \item 统计方法(n-gram, HMM)使语言处理变得可计算
    \item 神经网络方法(Word2Vec, Transformer)让机器自动学习语言结构
    \item 大语言模型(LLMs)在几乎全部语言任务上达到接近人类的水平
\end{itemize}

语言的描写从定性走向定量，从手工规则走向数据驱动。形式化程度不断提高，而语言学理论(类型学、形态音位学、依存语法)也在现代NLP系统中找到了落地的形态。语言学至此完成了一次螺旋上升：古代的形式化直觉、19世纪的经验实证、20世纪的结构与生成理论、21世纪的计算实现。

\subsection{本章小结}

语言学的发展经历了"形式化直觉(古代)→ 系统描写(19世纪)→ 结构分析(20世纪初)→ 生成理论(20世纪中)→ 认知与功能并重(20世纪后期)→ 计算实现(21世纪)"的历程。贯穿其中的主线，是对"语言规律能否被精确、形式化地刻画"这一问题的不断探索。理解了这段历史，也就理解了为什么语言学可以用数学公式、自动机理论和逻辑演算来表述——这正是本书的出发点。

\newpage
\section{语言的分类}

\subsection{分类的两条路径}

对人类语言进行分类，主要有两条互补的路径：
\begin{itemize}
    \item \textbf{谱系分类(Genetic classification)}: 依据语言的\emph{历史亲缘关系}，把有共同祖先的语言归入同一个语系。
    \item \textbf{类型学分类(Typological classification)}: 依据语言的\emph{结构特征}，把语法特征相似的语言归为一类，不论它们是否有亲缘关系。
\end{itemize}

这两条路径回答了不同的问题：谱系分类回答"语言从哪里来"，类型学分类回答"语言是什么样"。它们互补：同语系的语言通常共享类型特征(如印欧语多属屈折语)，但类型上相似的语言可以分属不同语系(如汉语、越南语、泰语都是孤立语却分属不同语系)。

\subsection{谱系分类}

\subsubsection{比较法与同源词}

谱系分类的基础是\emph{历史比较语言学}。通过比较不同语言中的同源词(cognates)，发现系统性的语音对应规律，可以证明它们同源并重建祖语。同源词的判定标准是"系统的语音对应"，而不是表面相似。例如英语 \emph{father}、德语 \emph{Vater}、拉丁语 \emph{pater}、梵语 \emph{pit\=r} 表面差异很大，但呈现出系统的对应(p-f-p-p)，由此证明同源。

\begin{equation}
    \text{Proto-Language} \xrightarrow{\text{分化}} \{L_1, L_2, \ldots, L_n\}
\end{equation}

\subsubsection{谱系树模型与波浪模型}

谱系树模型(树状图)假设语言从一个祖语"树状"地分化。但语言接触会使分化后的语言重新趋同，因此历史语言学还提出了波浪模型(Wave Model)：语言变化像波浪一样从创新中心向四周扩散，相邻语言共享某些变化。现代计算语言学用系统发生树(Phylogenetic tree)方法，基于词汇、语音、语法数据自动重建语言家族树。

\subsubsection{世界主要语系}

\begin{table}[H]
\centering
\caption{世界主要语系(使用人口为约数)}
\begin{tabulary}{\textwidth}{@{}CCC@{}}
\toprule
\textbf{语系} & \textbf{代表语言} & \textbf{使用人口(亿)} \\
\midrule
印欧语系 & 英语、西班牙语、印地语、俄语 & 32 \\
汉藏语系 & 汉语、藏语、缅甸语 & 14 \\
尼日尔-刚果语系 & 斯瓦希里语、约鲁巴语 & 7 \\
亚非语系 & 阿拉伯语、希伯来语、豪萨语 & 5 \\
南岛语系 & 马来语、他加禄语、马达加斯加语 & 3.8 \\
达罗毗荼语系 & 泰米尔语、泰卢固语 & 2.5 \\
突厥语系 & 土耳其语、哈萨克语、维吾尔语 & 2 \\
\bottomrule
\end{tabulary}
\end{table}

\subsubsection{语系的层次结构}

语系的内部有层次结构，从大到小依次为：语系(Family)→ 语族(Subfamily)→ 语支(Branch)→ 语言(Language)→ 方言(Dialect)。以印欧语系为例：
\begin{itemize}
    \item \textbf{日耳曼语族}: 西日耳曼语支(英语、德语、荷兰语)、北日耳曼语支(瑞典语、挪威语、丹麦语)
    \item \textbf{罗曼语族}: 法语、西班牙语、意大利语、葡萄牙语、罗马尼亚语
    \item \textbf{斯拉夫语族}: 俄语、波兰语、捷克语、塞尔维亚语
    \item \textbf{印度-伊朗语族}: 印地语、乌尔都语、波斯语、孟加拉语
\end{itemize}

值得注意的是，语言与方言的界限往往是社会政治的而非纯语言学的："一种语言就是拥有一支军队的方言"这一说法形象地说明了这一点。汉语的各个"方言"(如粤语、闽南语、吴语)彼此不能通话，差异程度堪比欧洲的不同语言，但因共享统一的文字与民族认同而被归为"汉语的方言"。

\subsection{类型学分类:按词法类型}

语言按构词方式可分为四大类型，这是最经典的类型学分类：

\subsubsection{孤立语(Isolating)}

孤立语中，词几乎没有形态变化，语法关系主要靠\emph{语序}和\emph{虚词}表达。词、语素、音节之间基本一一对应：
\begin{equation}
    \text{例: 汉语 "我吃饭"} \quad \text{每个字 = 一个语素 = 一个词}
\end{equation}
汉语、越南语、泰语、约鲁巴语都是典型的孤立语。孤立语没有"复数词缀""格词缀"这类屈折变化："我吃饭"和"他吃饭"中"吃"不因主语的人称而变形。

\subsubsection{黏着语(Agglutinative)}

黏着语的词干后逐个添加词缀，每个词缀只表达一种语法信息，词缀边界清晰，可以机械地堆叠：
\begin{equation}
    \text{例: 土耳其语 ev-ler-im-den} = \text{房子-复数-我的-从}
\end{equation}
\textit{ev-ler-im-den} 可以拆解为词干 \textit{ev}(房子)、复数词缀 \textit{-ler}、第一人称属有词缀 \textit{-im}(我的)、离格词缀 \textit{-den}(从)。每个词缀只有一个意义，边界一目了然。日语、韩语、土耳其语、斯瓦希里语、匈牙利语都是黏着语。

\subsubsection{屈折语(Fusional)}

屈折语中，一个词缀可能同时携带多种语法信息，词缀与意义之间是多对多的融合关系：
\begin{equation}
    \text{例: 拉丁语 amo} = \text{我爱(一个词尾 -o 同时表示第一人称、单数、现在时、陈述语气)}
\end{equation}
屈折语的形态与语序有"冗余"：拉丁语中 \textit{puella amat puerum}(女孩爱男孩)即使换语序也基本可解，因为格词尾标明了谁爱谁。俄语、德语、西班牙语、阿拉伯语、古典希腊语都是典型的屈折语。

\subsubsection{多式综合语(Polysynthetic)}

多式综合语(也称复综语)的一个动词词中可以包含主语、宾语、副词、时态等几乎所有成分，一个词就可以表达一个完整的句子：
\begin{equation}
    \text{例: 因纽特语 tusaatsiarunnanngittualuujunga} = \text{我听不太清}
\end{equation}
北美原住民语言(如莫霍克语)、西伯利亚的楚科奇语、澳大利亚的许多原住民语言都是多式综合语。

\begin{table}[H]
\centering
\caption{形态类型对比}
\begin{tabulary}{\textwidth}{@{}CCCCC@{}}
\toprule
\textbf{特征} & \textbf{孤立语} & \textbf{黏着语} & \textbf{屈折语} & \textbf{多式综合语} \\
\midrule
语素/词比 & 约1:1 & 高 & 中 & 极高 \\
词缀边界 & 无词缀 & 清晰 & 融合 & 高度融合 \\
代表语言 & 汉语、越南语 & 日语、土耳其语 & 德语、俄语 & 因纽特语、莫霍克语 \\
语法手段 & 语序、虚词 & 词缀堆叠 & 词尾变化 & 词内整合 \\
\bottomrule
\end{tabulary}
\end{table}

需要强调的是，这四类只是\emph{理想类型}。真实语言几乎都混合了多种类型：英语在形态上是孤立的(几乎没有格变化)，但保留了屈折遗迹(\emph{-s}复数、\emph{-ed}过去时、不规则动词 \emph{go/went})。因此类型学分类描述的是语言的\emph{主导倾向}，而不是非此即彼的标签。

\subsection{类型学分类:按句法类型}

\subsubsection{基本语序}

按句中主语(S)、动词(V)、宾语(O)的相对顺序，世界语言可分为若干类型：

\begin{table}[H]
\centering
\caption{基本语序类型及其分布(基于WALS数据库)}
\begin{tabulary}{\textwidth}{@{}CCC@{}}
\toprule
\textbf{语序} & \textbf{比例} & \textbf{代表语言} \\
\midrule
SVO & 42\% & 英语、汉语、法语、俄语 \\
SOV & 45\% & 日语、土耳其语、藏语、印地语 \\
VSO & 9\% & 阿拉伯语、威尔士语、爱尔兰语 \\
VOS & 3\% & 马尔加什语、菲吉语 \\
OVS / OSV & $<1\%$ & 罕见(如希克斯卡亚纳语) \\
\bottomrule
\end{tabulary}
\end{table}

值得注意的是，SOV(主语-宾语-动词)和SVO(主语-动词-宾语)合计占87\%，而OVS/OSV几乎不存在。这不是巧合，而是反映了语言对人类最自然的信息组织方式(施事在前)。

\subsubsection{类型学相关性:头参数}

语序不是孤立的，它与介词类型、名词-形容词顺序等"打包"出现，形成类型学上的相关性(相关词序共性)。核心规律是\emph{中心词(head)参数}：
\begin{itemize}
    \item \textbf{中心词在前(Head-initial)}: 使用前置词(\textit{in the house})，SVO语序
    \item \textbf{中心词在后(Head-final)}: 使用后置词(\textit{家の中} = 家之中)，SOV语序
\end{itemize}

\begin{equation}
    \forall L: \text{Head-initial}(L) \leftrightarrow \text{Prep}(L)
\end{equation}

日语"家の中"(house-of-in)中，属格标记"の"位于名词之后，是典型的中心词在后(后置词)类型；英语 \textit{in the house} 则反之。汉语是少见的例外：基本语序为SVO(中心词在前)，却使用后置词式的结构("桌子上")，同时拥有前置词("在桌上")。这种"混合"说明类型学相关是统计规律而非铁律。

\subsection{形式化类型学特征}

语言类型学可以形式化地表述为\emph{语言共性}(Universals)。格林伯格(Greenberg)提出了大量蕴含式共性，可用一阶逻辑表示：

\begin{equation}
    \forall L \in \text{Languages}: \text{VSO}(L) \rightarrow \text{Prep}(L)
\end{equation}

格林伯格的经典共性包括：
\begin{itemize}
    \item 若语言以VSO为基本语序，则使用前置词(格林伯格共性3)
    \item 若语言以SOV为基本语序，则几乎所有属格都位于名词之前
    \item 若语言是主格-宾格语言，则宾语在动词之后或之前与语序相关
\end{itemize}

形式化的共性刻画允许我们定义类型空间：
\begin{equation}
    \text{语言类型} = \prod_{i} \{\text{特征 } F_i \text{ 的取值}\}
\end{equation}
每一种自然语言都可以表示成一组特征取值构成的"向量"，类型学的研究就是在该空间中寻找聚类与规律。

\subsubsection{蕴含层级与标记性}

格林伯格共性常常构成\emph{蕴含层级}(implicational hierarchy)。以格系统为例：
\begin{equation}
    \text{主格} \prec \text{宾格/通格} \prec \text{与格} \prec \text{属格} \prec \text{方位格} \prec \ldots
\end{equation}
层级意味着：一种语言若具有层级中靠后的格，则必然也具有其前面的所有格。这种蕴含关系与\emph{标记性}(markedness)理论一致：层级靠后的格是"有标记"(更特殊)的范畴。

\subsection{语言接触与混合}

\subsubsection{语言联盟(Sprachbund)}

地理上相邻但无亲缘关系的语言，长期接触后会在结构上趋同，形成语言联盟。最著名的例子是\emph{巴尔干语言联盟}：保加利亚语、罗马尼亚语、阿尔巴尼亚语、马其顿语分属不同语族，却共享"后置定冠词""复合将来时""属格/与格合并"等特征。语言联盟可以形式化为：
\begin{equation}
    L_1, L_2 \xrightarrow{\text{长期接触}} \text{共享特征集合 } F
\end{equation}

汉语与周边语言的接触(如日语、朝鲜语、越南语借入大量汉语词汇，粤语、闽南语向东南亚语言传播)是语言接触影响类型面貌的又一例证。

\subsubsection{皮钦语与克里奥尔语}

\begin{itemize}
    \item \textbf{皮钦语(Pidgin)}: 在没有共同语言的人群交流中产生的简化混合语，语法简单、词汇有限，\emph{没有母语者}。
    \item \textbf{克里奥尔语(Creole)}: 皮钦语成为下一代儿童的母语后，语法迅速复杂化，形成完整的自然语言。
\end{itemize}

\begin{equation}
    \text{Pidgin}(L_1, L_2) \xrightarrow{\text{母语化}} \text{Creole}
\end{equation}

克里奥尔语的形成过程是"语言的出生"，在语言学上意义重大：它证明人类儿童在输入极不充分的条件下(皮钦语语法高度简化)仍能生成出语法完整的语言——这成为刺激贫乏论与普遍语法的天然证据。海地克里奥尔语(以法语为基础)是有数百万使用者的克里奥尔语。

\subsection{语言演化模型}

语言的演化和生物进化类似，可以用数学模型刻画。语言特征的使用比例随世代变化，最简单的模型是逻辑斯蒂增长模型：
\begin{equation}
    \frac{dP}{dt} = \mu \cdot P \cdot (1 - P)
\end{equation}
其中 $P$ 是某语言特征的使用比例，$\mu$ 是变化速率参数。

更完整的模型把语言变化看作\emph{复制者动态}(Replicator dynamics)：说话者学习语言时以一定概率复制母语者的特征，同时存在创新与选择压力。词汇替代率可以用斯瓦德什核心词表(Swadesh list)和词汇统计年代学(Glottochronology)来估计：
\begin{equation}
    t = \frac{\ln (c/C)}{\ln (\lambda)}
\end{equation}
其中 $c$ 是同源词保留比例，$C$ 是初始比例，$\lambda$ 是词汇保留率常数。现代计算类型学则直接用贝叶斯系统发生学方法，在语系树与时间尺度上同时估计音变、词变与类型演化的速率。

\subsection{本章小结}

语言分类的两条路径——谱系分类(历史亲缘)与类型学分类(结构特征)——分别从"源流"与"面貌"两个维度刻画语言。语系揭示语言的血缘关系，类型揭示语言的内部结构规律。类型学特征不是孤立的，而是以蕴含层级、中心词参数等方式相互关联。这些关联为最后一章"不同语言之间的差异"提供了系统的坐标：任何两种语言，都可以在语音、词法、句法、文字等维度上定位其差异。语言接触则提醒我们，语言不是封闭系统，"差异"本身也在不断演变。

\newpage
\section{形式化语言基础}

\subsection{形式语言理论概述}

形式语言理论是计算机科学和语言学的交叉领域，由诺姆·乔姆斯基在1956年奠基。它提供了一套精确的数学工具来描述语言的\emph{句法}结构。这里所说的"语言"不仅指自然语言，也包括编程语言、正则表达式、DNA序列等一切由符号构成的集合。

形式语言理论的核心思想：一个语言不过是某个字母表上的字符串集合，而文法(grammar)是生成这个集合的有限规则系统。它把"什么是合法的句子"这个语言学问题转化为"集合论与自动机"的数学问题：
\begin{equation}
    \text{语言} \subseteq \Sigma^* \quad \text{其中 } \Sigma \text{ 是字母表}
\end{equation}

\subsection{字母表与字符串}

\subsubsection{基本定义}

\begin{definition}[字母表]
字母表 $\Sigma$ 是一个有限的非空符号集合。
\end{definition}

\begin{example}
\begin{align*}
    \Sigma_{\text{英语}} &= \{a, b, c, \ldots, z\} \\
    \Sigma_{\text{二进制}} &= \{0, 1\} \\
    \Sigma_{\text{DNA}} &= \{A, C, G, T\}
\end{align*}
\end{example}

\begin{definition}[字符串]
字符串 $w$ 是字母表 $\Sigma$ 中符号的有限序列。
\end{definition}

\subsubsection{字符串的基本运算}

\begin{itemize}
    \item $|w|$: 字符串 $w$ 的长度
    \item $\epsilon$: 空字符串，$|\epsilon| = 0$
    \item $uv$: 字符串 $u$ 和 $v$ 的\emph{连接}(concatenation)
    \item $w^R$: 字符串 $w$ 的\emph{反转}(reverse)
    \item $w^k$: 字符串 $w$ 的 $k$ 次重复
    \item $\Sigma^*$: $\Sigma$ 上所有字符串的集合(包括 $\epsilon$)
    \item $\Sigma^+$: $\Sigma$ 上所有非空字符串的集合
\end{itemize}

\subsubsection{语言作为集合}

\begin{definition}[语言]
$\Sigma$ 上的一个语言 $L$ 是 $\Sigma^*$ 的任意子集：$L \subseteq \Sigma^*$。
\end{definition}

语言可以非常小(空集 $\emptyset$、单个字符串 $\{ab\}$)也可以无限大(如"所有以 $a$ 开头、$b$ 结尾的字符串" $\{awb : w \in \Sigma^*\}$)。关键问题是：\emph{一个无限语言能否用有限的方式描述?} 这正是文法和自动机要回答的问题。

语言之间有基本的集合运算：并、交、差、补、连接、Kleene星：
\begin{equation}
    L_1 \cdot L_2 = \{ uv \mid u \in L_1, v \in L_2 \}, \qquad L^* = \bigcup_{k \geq 0} L^k
\end{equation}

\subsection{文法与乔姆斯基层级}

\subsubsection{形式文法定义}

\begin{definition}[短语结构文法]
文法 $G$ 是一个四元组：
\begin{equation}
    G = \langle V_N, V_T, P, S \rangle
\end{equation}
其中：
\begin{itemize}
    \item $V_N$: 非终结符(Non-terminal)集合，代表语法范畴(如 S、NP、VP)
    \item $V_T$: 终结符(Terminal)集合，代表实际词/符号，$V_N \cap V_T = \emptyset$
    \item $P$: 产生式(Production)规则集合，形如 $\alpha \rightarrow \beta$
    \item $S \in V_N$: 起始符号(Start symbol)
\end{itemize}
\end{definition}

文法生成语言的过程是：从起始符 $S$ 出发，反复应用产生式，直到所有非终结符都被替换为终结符。记 $\Rightarrow$ 为"一步推导"，则
\begin{equation}
    L(G) = \{ w \in V_T^* \mid S \Rightarrow^* w \}
\end{equation}
其中 $\Rightarrow^*$ 表示零步或多步推导。一个语言若存在某个文法生成它，就称该语言是\emph{递归可枚举}的(这里指0型)。

\subsubsection{乔姆斯基层级}

根据产生式形式的限制，文法可以分为四个层次，每一层生成的语言严格包含下一层：

\begin{table}[H]
\centering
\caption{乔姆斯基层级}
\begin{tabulary}{\textwidth}{@{}CCCC@{}}
\toprule
\textbf{类型} & \textbf{文法} & \textbf{产生式形式} & \textbf{自动机} \\
\midrule
0 & 无限制文法 & $\alpha \rightarrow \beta$ & 图灵机 \\
1 & 上下文有关 & $\alpha A \beta \rightarrow \alpha \gamma \beta$ & 线性有界自动机 \\
2 & 上下文无关 & $A \rightarrow \gamma$ & 下推自动机 \\
3 & 正则 & $A \rightarrow aB$ 或 $A \rightarrow a$ & 有限状态自动机 \\
\bottomrule
\end{tabulary}
\end{table}

\begin{equation}
    \mathcal{L}_0 \supsetneq \mathcal{L}_1 \supsetneq \mathcal{L}_2 \supsetneq \mathcal{L}_3
\end{equation}
其中 $\mathcal{L}_i$ 表示 $i$ 型文法生成的语言类。层级关系是严格的：$\{a^n b^n\}$ 是上下文无关但非正则的；$\{a^n b^n c^n\}$ 是上下文有关但非上下文无关的。

\subsection{正则语言(Type-3)}

\subsubsection{正则表达式}

正则表达式(Regular Expression)是最直观地描述正则语言的方式，其递归定义如下：
\begin{equation}
    \begin{aligned}
    &\emptyset \text{ 是正则表达式(表示空语言)} \\
    &\epsilon \text{ 是正则表达式(表示 } \{\epsilon\} \text{)} \\
    &a \in \Sigma \text{ 是正则表达式(表示 } \{a\} \text{)} \\
    &\text{若 } r, s \text{ 是正则表达式，则 } (r|s), (rs), (r^*) \text{ 也是}
    \end{aligned}
\end{equation}
这里 $r|s$ 是并、$rs$ 是连接、$r^*$ 是Kleene星(零次或多次重复)。语法糖包括：$r^+ = rr^*$(一次以上)、$r? = (r|\epsilon)$(零次或一次)、$[a-z]$ 表示区间。

\subsubsection{有限状态自动机(FSA)}

\begin{definition}[DFA]
确定性有限自动机 $M = \langle Q, \Sigma, \delta, q_0, F \rangle$
\begin{itemize}
    \item $Q$: 有限状态集合
    \item $\Sigma$: 输入字母表
    \item $\delta: Q \times \Sigma \rightarrow Q$: 转移函数
    \item $q_0 \in Q$: 起始状态
    \item $F \subseteq Q$: 接受状态集合
\end{itemize}
\end{definition}

DFA从 $q_0$ 出发，逐个读入输入符号并转移状态。若读完整个输入后停在 $F$ 中的某个状态，则接受该字符串；否则拒绝。DFA接受的语言记作 $L(M)$。

\begin{theorem}[克莱因定理]
正则语言类 = 正则表达式可描述的语言类 = DFA/NFA可接受的语言类。
\end{theorem}

\begin{example}[识别英语复数后缀]
\begin{equation}
    \text{正则表达式: } [a-z]^* (s|es|ies)
\end{equation}
DFA能识别"字符串由小写字母组成且以 s、es 或 ies 结尾"。有限状态自动机的局限也在于此：它\emph{没有记忆}，无法统计括号匹配的嵌套层数，因此不能识别形如 $a^n b^n$ 的语言。
\end{example}

正则语言的封闭性质使其在计算上极其便利：正则语言对并、交、差、补、反转、同态都是封闭的，且可判定性(成员问题、空性、等价性)都有高效算法。这就是为什么正则表达式被广泛应用于文本搜索、词法分析(如编译器的lex)和NLP的词形还原。

\subsection{上下文无关文法(Type-2)}

\subsubsection{CFG定义}

上下文无关文法的产生式形式为：
\begin{equation}
    A \rightarrow \alpha \quad \text{其中 } A \in V_N, \alpha \in (V_N \cup V_T)^*
\end{equation}
"上下文无关"的意思是：替换一个非终结符时，不依赖于它左右两边的上下文。这使CFG非常适合描述语言的\emph{递归嵌套结构}——自然语言和编程语言中最常见的结构。

\begin{example}[回文语言]
$L = \{ w \in \{a,b\}^* \mid w = w^R \}$ 是上下文无关的：
\begin{equation}
    S \rightarrow aSa \mid bSb \mid a \mid b \mid \epsilon
\end{equation}
这个语言不是正则的：识别它需要记住已经读入的内容，而DFA没有记忆。
\end{example}

\subsubsection{句法树}

CFG的重要性质是：一个推导过程对应一棵\emph{句法树}(Parse Tree，也称推导树)，树根是 $S$，叶节点是终结符。句法树显式刻画了句子的层次结构，这正是自然语言句法分析的目标：

\begin{forest}
for tree={s sep=10mm, l sep=10mm}
[S
    [NP
        [Det [The]]
        [N [cat]]
    ]
    [VP
        [V [sat]]
        [PP
            [P [on]]
            [NP
                [Det [the]]
                [N [mat]]
            ]
        ]
    ]
]
\end{forest}

\subsubsection{歧义与二义性文法}

同一个句子可以对应多棵句法树，称为\emph{结构歧义}(structural ambiguity)。经典例子：
\begin{equation}
    \text{"I saw the man with the telescope"}
\end{equation}
可以是"我用望远镜看那个人"($with$ 修饰 $saw$)，也可以是"我看见那个拿望远镜的人"($with$ 修饰 $man$)。两种解读对应两棵不同的句法树。一个文法若对某个句子能产生多棵句法树，就称为二义性文法。存在无法消除二义性的上下文无关语言，即固有歧义语言。

\subsubsection{范式}

\begin{itemize}
    \item \textbf{乔姆斯基范式(CNF)}: 每条产生式形如 $A \rightarrow BC$ 或 $A \rightarrow a$。任何CFG都可以等量地(语言不变)转化为CNF，且任何不产生 $\epsilon$ 的CFG都可以。CNF的最大价值在于：它把推导树的内部节点统一为"两分叉"，使得\emph{CYK算法}(动态规划句法分析)可以直接应用。
    \item \textbf{格雷巴赫范式(GNF)}: 产生式形如 $A \rightarrow a\alpha$，其中 $\alpha \in V_N^*$。GNF保证每一步都吃掉一个终结符，便于构造下推自动机。
\end{itemize}

\subsection{下推自动机与上下文无关语言}

上下文无关语言对应\emph{下推自动机}(PDA)：在有限状态的基础上增加一个栈(stack)。栈提供了"无限记忆"，但只能以\emph{后进先出}方式访问。正是这个栈使得PDA能识别 $a^n b^n$(读 $a$ 时压栈，读 $b$ 时弹栈)。

\begin{theorem}
$L$ 是上下文无关语言当且仅当存在一个下推自动机 $M$ 使得 $L = L(M)$。
\end{theorem}

\subsection{上下文有关文法(Type-1)}

产生式形式：
\begin{equation}
    \alpha A \beta \rightarrow \alpha \gamma \beta
\end{equation}
其中 $A \in V_N$，$\alpha, \beta, \gamma \in (V_N \cup V_T)^*$，且 $\gamma \neq \epsilon$。

这里的限制是：$A$ 只有在左右分别是 $\alpha$ 和 $\beta$ 时才能被替换为 $\gamma$，即替换依赖上下文。这种"上下文敏感性"使上下文有关文法能表达非上下文无关的结构。

\begin{example}
语言 $\{a^n b^n c^n : n \geq 1\}$ 是上下文有关的。上下文无关文法无法识别它，因为 $a$、$b$、$c$ 的个数需要同步约束；而上下文有关文法可以让每个 $a$ 在与某个 $b$、某个 $c$ 配对后一起替换：
\begin{equation}
    S \rightarrow aSBC \mid aBC, \qquad CB \rightarrow BC, \qquad aB \rightarrow ab, \quad bB \rightarrow bb, \quad bC \rightarrow bc, \quad cC \rightarrow cc
\end{equation}
上下文有关文法对应\emph{线性有界自动机}(LBA)，即使用量受限的图灵机(磁带长度与输入长度成正比)。
\end{example}

\subsection{形式语言与自然语言}

\subsubsection{自然语言的复杂度}

自然语言的句法结构位于乔姆斯基层级的哪一层？这是形式语言学的基本问题之一。答案是：自然语言超出上下文无关文法的能力，但通常只需\emph{温和上下文有关}(Mildly context-sensitive)的层级：

\begin{itemize}
    \item \textbf{跨串依赖(Cross-serial dependencies)}: 瑞士德语和荷兰语中的交叉依赖结构。在"我让他叫约翰帮忙打扫房间"这类结构中，动词与宾语交叉对应，无法用CFG表达。
    \item \textbf{温和上下文有关(Mildly context-sensitive)}: 介于CFG与上下文有关文法之间的层级，能表达交叉依赖但保持可解析性。
    \item \textbf{树邻接文法(TAG)}、\textbf{组合范畴文法(CCG)}: 是温和上下文有关的代表性形式体系，被广泛用于计算句法学。
\end{itemize}

\begin{equation}
    \mathcal{L}_{\text{CFG}} \subsetneq \mathcal{L}_{\text{TAG}} \subsetneq \mathcal{L}_{\text{MCS}} \subseteq \mathcal{L}_1
\end{equation}

\subsubsection{依存语法}

短语结构文法描述的是"成分层级"，而依存语法(Dependency Grammar)描述\emph{词与词之间的二元关系}：
\begin{equation}
    \text{依存关系} = \{(h, d, r) \mid h, d \in W, r \in R\}
\end{equation}
其中 $h$ 是中心词，$d$ 是从属词，$r$ 是关系类型(如 nsubj、dobj)。依存结构更接近语义关系，且与语序相对独立，因此跨语言的适用性很强——这也为最后一章"不同语言的差异"提供了统一的比较工具：不同语言在依存关系上高度相似，差异主要体现在语序和形态。

\subsection{形式语义基础}

\subsubsection{Lambda演算}

形式语言理论主要关心句法(什么是合法句子)，而\emph{形式语义学}要回答"句子是什么意思"。Lambda演算提供了组合意义的形式工具：
\begin{equation}
    \begin{aligned}
    &\text{变量: } x, y, z, \ldots \\
    &\text{抽象: } \lambda x.M \quad (\text{函数，把 } x \text{ 映射到 } M) \\
    &\text{应用: } (M \ N) \quad (\text{把函数 } M \text{ 应用到 } N)
    \end{aligned}
\end{equation}

\begin{example}
"runs" 的语义表示：
\begin{equation}
    \llbracket \text{runs} \rrbracket = \lambda x.\text{run}(x)
\end{equation}
它表示"谓词 run 对主体 x 成立"这个函数。组合语义的基本思想是：一个句子的意义由各成分的意义\emph{组合}而成，Lambda演算使这种组合可以机械地进行。
\end{example}

\subsection{本章小结}

形式语言理论为语言学提供了精确的"标尺"：乔姆斯基层级刻画了语言结构复杂度的层级，正则语言对应最"简单"的结构(无递归记忆)，上下文无关文法对应嵌套结构(自然语言句法的核心层级)，温和上下文有关覆盖了自然语言的完整复杂度。文法则架起了"有限规则"与"无限语言"之间的桥梁——这正回答了语言习得中的核心悖论：人类如何用有限的头脑掌握无限的语言能力。下一章将从词出发，考察词的内部结构(形态学)，那里同样充满形式化的规律。

\newpage
\section{词法(形态学)}

\subsection{形态学概述}

形态学(Morphology)研究词的内部结构和构词规则。词不是不可分割的意义单元——大多数词都可以分解为更小的、携带意义的构件。

\begin{definition}[语素]
语素(Morpheme)是最小的有意义的语言单位，不能再继续分割为更小的有意义的部件。
\end{definition}

例如 \emph{unhappiness} 可以分解为 $un$ + $happi$ + $ness$ 三个语素，而 \emph{cat} 虽然可分割为语音 /kæt/，但那些音段本身没有意义，因此 \emph{cat} 是一个语素。

形态学研究的两大任务：
\begin{itemize}
    \item \textbf{屈折(Inflection)}: 同一个"词"的不同语法形式(如 walk/walked/walking)
    \item \textbf{派生(Derivation)}: 由已有词构成新词(happy → happiness)
\end{itemize}

形态学还提供了一个重要的类型学参数：语素与词的比率——这个比率决定了语言的形态类型(孤立语/黏着语/屈折语/多式综合语)，是最后一章跨语言比较的核心维度之一。

\subsection{语素类型}

\subsubsection{按独立性}

\begin{itemize}
    \item \textbf{自由语素(Free morpheme)}: 可以独立成词，单独使用
    \begin{equation}
        \text{例: cat, run, happy, 桌子}
    \end{equation}
    
    \item \textbf{黏着语素(Bound morpheme)}: 必须依附于其他语素，不能单独成词
    \begin{equation}
        \text{例: -s, -ed, un-, -ness, 汉语中的 "们" 在单独出现时无意义}
    \end{equation}
\end{itemize}

\subsubsection{按功能}

\begin{itemize}
    \item \textbf{词根(Root)}: 词的核心意义，去掉所有词缀后剩下的部分
    \item \textbf{词干(Stem)}: 可以添加屈折词缀的形式(可能是"词根+派生词缀")
    \item \textbf{词缀(Affix)}: 附着在词根/词干上的黏着语素，按位置分类：
    \begin{itemize}
        \item 前缀(Prefix): un-, re-, pre- (位于词根之前)
        \item 后缀(Suffix): -s, -ed, -ing, -ness (位于词根之后，世界上最多见)
        \item 中缀(Infix): 插入词根内部，较少见，如他加禄语 \textit{-um-}
        \item 环缀(Circumfix): 同时加在词根前后，如德语过去分词 ge-mach-t
        \item 异干法(Suppletion): 完全不同形式的语素交替，如 go/went
    \end{itemize}
\end{itemize}

\subsection{屈折(Inflection)}

\subsubsection{屈折的性质}

屈折不改变词的词汇意义和词类，只添加\emph{语法}信息(如数、格、时态)，并且一个屈折范畴是强制性的——英语名词一旦用作复数就必须加 \emph{-s}，不能省略：
\begin{equation}
    \text{walk} \xrightarrow{+\text{过去时}} \text{walked}
\end{equation}

\subsubsection{屈折范畴}

\begin{itemize}
    \item 数(Number): cat/cats，汉语"猫/猫们"(有限制，只用于人)
    \item 格(Case): he/him/his —— 英语保留了主格/宾格/属格残迹
    \item 时态(Tense): walk/walked —— 时间定位
    \item 体(Aspect): walk/walking —— 动作的内部时间结构(进行/完成)
    \item 人称(Person): am/are/is
    \item 性(Gender): 德语 der/die/das，法语 le/la —— 语法性不一定是自然性别
    \item 语气(Mood): 陈述/虚拟/祈使 (英语 subjunctive 残存于 "I suggest that he \emph{be} here")
    \item 态(Voice): 主动/被动 (was eaten)
    \item 级(Degree): tall/taller/tallest (形容词屈折)
\end{itemize}

\subsubsection{屈折的跨语言差异}

屈折范畴在不同语言中的分布差异极大，这是跨语言比较的重点之一：
\begin{itemize}
    \item 汉语几乎无屈折，时态靠时间词和体助词("了""着""过")表达
    \item 英语保留了少量屈折(复数-s、过去时-ed、第三人称-s)，但比德语、俄语少得多
    \item 俄语名词有6个格，形容词有24种形式(数×性×格)
    \item 日语、韩语动词有人称无关、但极复杂的敬体/简体系统
    \item 许多语言名词没有"格"，但有的语言有几十种格(如芬兰语名词有15个格)
\end{itemize}

\subsection{派生(Derivation)}

\subsubsection{派生的性质}

派生通过加词缀创造\emph{新词}，可以改变词类或改变核心意义，且派生是选择性的(不一定能派生所有同类词)：
\begin{align}
    \text{happy (adj.)} &\xrightarrow{+\text{-ness}} \text{happiness (n.)} \\
    \text{write (v.)} &\xrightarrow{+\text{re-}} \text{rewrite (v.)} \\
    \text{kind (adj.)} &\xrightarrow{+\text{un-}} \text{unkind (adj.)} \\
    \text{teach (v.)} &\xrightarrow{+\text{-er}} \text{teacher (n.)}
\end{align}

\subsubsection{派生与屈折的对比}

\begin{table}[H]
\centering
\caption{屈折与派生的区别}
\begin{tabulary}{\textwidth}{@{}CCC@{}}
\toprule
\textbf{特征} & \textbf{屈折} & \textbf{派生} \\
\midrule
是否创造新词 & 否 & 是 \\
是否改变词类 & 否 & 常改变 \\
强制性 & 是 & 否 \\
与语法环境相关 & 是 & 否 \\
产出顺序 & 在派生之后 & 在屈折之前 \\
\bottomrule
\end{tabulary}
\end{table}

这一顺序(派生在内、屈折在外)是语言普遍存在的规律，可形式化为词缀顺序约束(见下文)。

\subsection{构词方式总览}

构词方式(Word Formation)研究新词是如何被创造出来的。除了前面介绍的屈折与派生，语言学家还归纳出大量其他构词方式，大致可分为"词内新造"与"外来借入"两大类，词内新造又分为组合型、缩短型与转换/语义型。下表给出全貌：

\begin{table}[H]
\centering
\caption{构词方式一览}
\begin{tabulary}{\textwidth}{@{}CCCC@{}}
\toprule
\textbf{类别} & \textbf{方式} & \textbf{机制} & \textbf{例} \\
\midrule
组合型 & 复合 & 词根+词根 & 火车、blackboard \\
组合型 & 重叠 & 重复语素 & 人人、anak-anak \\
组合型 & 拟声词 & 语音模拟声音 & buzz、叮当 \\
组合型 & 专名转化 & 专名转为通名 & sandwich、牛顿 \\
缩短型 & 截短 & 删去词的一部分 & lab、bus \\
缩短型 & 逆构词 & 删去误认的词缀 & edit$\leftarrow$editor \\
缩短型 & 缩写 & 取各词首字母 & radar、FBI \\
缩短型 & 混成 & 两词部分融合 & brunch、smog \\
转换/语义 & 转类 & 词形不变换词类 & email(v.) \\
转换/语义 & 语义造词 & 词义引申出新义 & web$\rightarrow$互联网 \\
转换/语义 & 比喻造词 & 隐喻/转喻造词 & 瓶颈、bug \\
借入 & 音译 & 借音不借形 & 沙发(sofa) \\
借入 & 意译 & 译义造新词 & 电脑(computer) \\
借入 & 仿译 & 逐语素对应翻译 & 热狗(hot dog) \\
借入 & 音义兼译 & 音义兼顾 & 可口可乐 \\
借入 & 借形 & 借字形不借音 & 电话(自日语) \\
借入 & 语义借用 & 旧词借新义项 & 病毒(计算机) \\
\bottomrule
\end{tabulary}
\end{table}

下面分组详述。

\subsection{组合型构词}

组合型构词把已有的词根或语素"拼"在一起，是汉语、英语、日语、德语中最能产的构词方式。

\subsubsection{复合(Compounding)}

复合(Compounding)是两个或以上的词根直接组合成新词的过程，产物叫复合词(compound)。它是世界上绝大多数语言都具备的构词方式。

\begin{table}[H]
\centering
\caption{复合词的内部结构类型(汉语为例)}
\begin{tabulary}{\textwidth}{@{}CCC@{}}
\toprule
\textbf{类型} & \textbf{结构关系} & \textbf{例子} \\
\midrule
并列复合 & 两成分地位平等 & 山水、黑白、声音 \\
偏正复合 & 前修饰、后中心 & 火车、黑板、白菜 \\
动宾复合 & 动词+宾语 & 开门、司令、起草 \\
动补复合 & 动词+补语 & 提高、证明、扩大 \\
主谓复合 & 主语+谓语 & 地震、心疼、民主 \\
\bottomrule
\end{tabulary}
\end{table}

英语名词复合词最常见的是"修饰语+中心语"结构(football = foot+ball，中心在"球")，但词性组合很多样：名词+名词(rainbow)、形容词+名词(blackboard)、动词+名词(swimming pool)等。

复合词与短语的区别在英语中可以用\emph{重音}来检验：复合词重音落在第一成分(blackboard 黑板)，而短语重音落在第二成分(black board 黑色的板)。汉语中复合词与短语的界限较模糊："火车""黑板"已经完全词汇化，而"大门""红车"既可看作复合词也可看作短语。

复合是汉语最能产、最古老的构词方式之一，现代汉语新词绝大多数是复合词：手机、高铁、网红、外卖。德语也以复合著称，可以造出很长的复合词(Rindfleischetikettierungsüberwachungsaufgabenübertragungsgesetz，牛肉标签监管任务转移法)；英语同样用复合不断造新词(firewall、lockdown、greenhouse)。

\subsubsection{重叠(Reduplication)}

重叠(Reduplication)通过重复语素或音节构成新词形，可以是全部重复，也可以是部分重复：

\begin{table}[H]
\centering
\caption{重叠的类型与功能}
\begin{tabulary}{\textwidth}{@{}CCC@{}}
\toprule
\textbf{类型} & \textbf{方式} & \textbf{例子} \\
\midrule
完全重叠 & 整个词根重复 & 人人、日日、anak-anak(孩子们) \\
部分重叠 & 只重复部分音节 & 他加禄语、拉丁语完成时(cecini) \\
变韵重叠 & 重叠时元音变化 & chit-chat、ping-pong、zigzag \\
\bottomrule
\end{tabulary}
\end{table}

重叠的功能因语言而异，非常多样：
\begin{itemize}
    \item 名词：表示复数或"每一"——印尼语 anak-anak(孩子们)、汉语"人人"(每个人)、"家家"
    \item 动词：表示尝试或随意——汉语"看看""想想""试试"
    \item 形容词：表示程度加深或生动——汉语"高高""红红的""漂漂亮亮"
    \item 量词：表示"逐个"——汉语"一个个""一天天"
    \item 拟声拟态：摹写声音或状态——日语"わくわく"(心怦怦)、"ざあざあ"(雨声)
\end{itemize}

汉语、印尼语、马来语、日语是重叠极丰富的语言；英语重叠较少，主要用于拟声(ding-dong)和口语叠词(tip-top、hush-hush)。

\subsubsection{拟声词(Onomatopoeia)}

拟声词(Onomatopoeia)用语音模拟事物声音，是"以声造词"。

\begin{itemize}
    \item \textbf{直接拟声}: 模拟自然界或动物的声音。英语 buzz(嗡嗡)、meow(喵)、clap(啪)、cuckoo(布谷)，汉语"叮当""哗啦""喵""汪汪"。
    \item \textbf{拟声拟态词(ideophone)}: 不仅模拟声音，还摹写形状、状态、心情。日语拟声拟态词极其丰富，数量达数千个，如"わくわく"(心怦怦跳)、"きらきら"(闪闪发光)、"ぐっすり"(熟睡)；韩语、西非语言也有庞大的拟态词系统。
\end{itemize}

拟声词具有高度的\emph{语言特异性}：同一种声音在不同语言中"听感"截然不同，因为每种语言都用自己有限的音位系统去模拟：

\begin{table}[H]
\centering
\caption{同一声音在不同语言中的拟声差异}
\begin{tabulary}{\textwidth}{@{}CCCC@{}}
\toprule
\textbf{声音} & \textbf{英语} & \textbf{日语} & \textbf{汉语} \\
\midrule
狗叫 & woof / bark & ワンワン (wan-wan) & 汪汪 (wāngwāng) \\
鸡叫 & cock-a-doodle-doo & コケコッコー (kokekokkō) & 喔喔喔 (wōwōwō) \\
猫叫 & meow & ニャー (nyā) & 喵 (miāo) \\
\bottomrule
\end{tabulary}
\end{table}

\subsubsection{专名转化(Eponym)}

专名转化(Eponym)把专有名词——人名、地名、品牌名、文学人物——转用为普通名词：

\begin{table}[H]
\centering
\caption{专名转化的来源}
\begin{tabulary}{\textwidth}{@{}CCC@{}}
\toprule
\textbf{来源} & \textbf{机制} & \textbf{例子} \\
\midrule
人名 & 以发明者或人物命名 & sandwich(三明治)、volt(伏特)、牛顿、瓦特 \\
地名 & 以产地命名 & champagne(香槟)、denim(牛仔布)、china(瓷器) \\
品牌 & 品牌名泛化 & kleenex(纸巾)、xerox(复印)、百度一下 \\
神话文学 & 以文学人物命名 & quixotic(唐吉诃德式)、阿斗 \\
\bottomrule
\end{tabulary}
\end{table}

品牌名"普通化"(genericide)是专名转化最活跃的类型：品牌一旦成为某类商品的通称，商标就失去专有性。kleenex 原是面巾纸品牌，现泛指纸巾；xerox 原是复印机品牌，现作动词"复印"。汉语"百度"也正在经历同样的转化："百度一下"就是"搜索一下"。历史人物"阿斗"(扶不起的阿斗)、"雷锋"(助人为乐者)则是人物转义的典型。

\subsection{缩短型构词}

缩短型构词从已有的长词或短语中删减材料，形成更短的新形式，追求简便。

\subsubsection{截短(Clipping)}

截短(Clipping)删去原词的一部分，词义与词类不变，只是形式变短：

\begin{table}[H]
\centering
\caption{截短的部位}
\begin{tabulary}{\textwidth}{@{}CCC@{}}
\toprule
\textbf{类型} & \textbf{方式} & \textbf{例子} \\
\midrule
保留词首 & 删去词尾 & exam(examination)、lab(laboratory)、phone(telephone) \\
保留词尾 & 删去词首 & bus(omnibus)、phone(telephone) \\
保留中部 & 删去首尾 & flu(influenza)、fridge(refrigerator) \\
\bottomrule
\end{tabulary}
\end{table}

截短词通常带口语、亲切的色彩，如 \emph{prof}（professor）、\emph{gym}（gymnasium）、\emph{math}（mathematics）、\emph{tech}（technology）。

人名也常被截短作昵称，如 \emph{Jen}（Jennifer）、\emph{Pat}（Patricia）。有些截短词完全取代了原词，如 \emph{bus} 已很少有人知道它来自 \emph{omnibus}。

\subsubsection{逆构词(Back-formation)}

逆构词(Back-formation)通过删去一个被\emph{误认为}词缀的部分，逆向造出"更简单"的词。其关键在于词形的历史顺序：先有长形式，后逆推出短形式。

\begin{table}[H]
\centering
\caption{逆构词例子}
\begin{tabulary}{\textwidth}{@{}CCC@{}}
\toprule
\textbf{新词} & \textbf{来源} & \textbf{机制} \\
\midrule
edit(编辑) & editor & 误以为 -er 是"施事者"后缀 \\
swindle(诈骗) & swindler & 同上 \\
babysit(照看婴儿) & babysitter & 删去 -er \\
donate(捐赠) & donation & 误以为 -ion 是名词化后缀 \\
pea(豌豆) & pease & 误以为 -s 是复数后缀 \\
\bottomrule
\end{tabulary}
\end{table}

逆构词与截短的区别：截短只是缩短形式、不改变词类与构词结构(lab 仍是名词)；逆构词则\emph{改变词类或结构}(editor 名词$\rightarrow$edit 动词)。逆构词反映说者对构词的"再分析"——把不存在的词缀去掉，类比地推出新词，甚至 laser 也被戏谑地逆构为动词 to lase(发射激光)。

\subsubsection{缩写(Abbreviation)}

缩写从短语中取首字母(或首音节)形成新词，是科技与网络时代最活跃的构词方式：

\begin{table}[H]
\centering
\caption{缩写的类型}
\begin{tabulary}{\textwidth}{@{}CCCC@{}}
\toprule
\textbf{类型} & \textbf{读音} & \textbf{例子} & \textbf{来源} \\
\midrule
首字母缩略词(Acronym) & 按单词拼读 & radar、laser、NATO、scuba & 各词首字母 \\
首字母缩写(Initialism) & 按字母逐个读 & FBI、TV、UN、IBM & 各词首字母 \\
音节缩略/简称 & 取首音节 & 北外、央视、スマホ、東大 & 首音节/简称 \\
\bottomrule
\end{tabulary}
\end{table}

\begin{itemize}
    \item \textbf{Acronym(按词拼读)}: radar = RAdio Detection And Ranging、laser = Light Amplification by Stimulated Emission of Radiation、NATO = North Atlantic Treaty Organization、scuba = self-contained underwater breathing apparatus。
    \item \textbf{Initialism(按字母读)}: FBI、TV、UN、IBM、USA。
    \item \textbf{音节缩略/简称}: 汉语"北外"(北京外国语大学)、"央视"(中央电视台)；日语"スマホ"(スマートフォン，smartphone)、"東大"(東京大学)；俄语也有大量音节缩略词。
\end{itemize}

\subsubsection{混成(Blending)}

混成(Blending)把两个词的部分音节"嫁接"成一个新词，又称拼合词(portmanteau)。与复合不同，混成通常只保留源词的部分音节：

\begin{table}[H]
\centering
\caption{混成的结构}
\begin{tabulary}{\textwidth}{@{}CCC@{}}
\toprule
\textbf{结构} & \textbf{例子} & \textbf{来源} \\
\midrule
前词首+后词尾 & brunch(早午餐) & breakfast+lunch \\
前词首+后词首 & modem(调制解调器) & modulator+demodulator \\
重叠混成(音相近) & guesstimate(毛估)、spork(勺叉) & guess+estimate、spoon+fork \\
\bottomrule
\end{tabulary}
\end{table}

其他经典混成：smog = smoke+fog(烟雾)、motel = motor+hotel(汽车旅馆)、sitcom = situation+comedy(情景喜剧)、telethon = television+marathon(电视马拉松)。"portmanteau"一词来自刘易斯·卡罗尔在《爱丽丝镜中奇遇》中造的词 chortle = chuckle+snort。

混成在英语中最为活跃；汉语因每个字都是独立语素，混成较少见，口语缩合(如"高大上")更接近简称而非典型混成，日语中也较少见。

\subsection{词类转换与语义构词}

\subsubsection{转类(Conversion)}

转类(Conversion)不改变词形、只改变词类，因此又称"零派生"(zero derivation)。它是英语最活跃的构词方式之一——英语屈折很弱，词本身几乎不带词类标记：

\begin{table}[H]
\centering
\caption{转类的方向}
\begin{tabulary}{\textwidth}{@{}CCC@{}}
\toprule
\textbf{转换方向} & \textbf{例子} & \textbf{说明} \\
\midrule
名词$\rightarrow$动词 & email(v.)、google(v.)、access(v.) & 网络时代尤为典型 \\
动词$\rightarrow$名词 & walk(n.)、run(n.)、guess(n.) & 动作转指事件/结果 \\
形容词$\rightarrow$动词 & empty(v.)、clean(v.)、dry(v.) & 使然动词 \\
形容词$\rightarrow$名词 & the rich、the poor、original & 指人/指事物 \\
\bottomrule
\end{tabulary}
\end{table}

汉语的转类表现为"兼类"与"词类活用"：名词"雷锋"可临时用作形容词("他很雷锋")；"百度"由名词转作动词("百度一下")。传统语法学界长期争论"汉语到底有没有词类"，正是因为汉语词几乎没有形态标记，同一个词形可以出现在多种语法位置。

\subsubsection{语义造词:词义变化(Semantic Change)}

语义造词不改变词形，而是通过\emph{词义变化}让已有词获得新义。历史语义学归纳出词义变化的几种基本类型：

\begin{table}[H]
\centering
\caption{词义变化的类型}
\begin{tabulary}{\textwidth}{@{}CCC@{}}
\toprule
\textbf{类型} & \textbf{机制} & \textbf{例子} \\
\midrule
词义扩大 & 指称范围变大 & 汉语"江""河"(专指长江黄河$\rightarrow$泛指河流)；英语 holiday(圣日$\rightarrow$假日) \\
词义缩小 & 指称范围变小 & 汉语"禽"(鸟兽总称$\rightarrow$家禽)；英语 meat(食物$\rightarrow$肉)、deer(动物$\rightarrow$鹿) \\
词义转移 & 所指对象改变 & 汉语"走"(跑$\rightarrow$走)、"涕"(眼泪$\rightarrow$鼻涕)；英语 gay(快乐$\rightarrow$同性恋) \\
词义贬化 & 感情色彩变坏 & 英语 villain(村民$\rightarrow$坏人)、silly(幸福$\rightarrow$愚蠢) \\
词义褒化 & 感情色彩变好 & 英语 knight(男孩$\rightarrow$骑士)、nice(愚蠢$\rightarrow$好) \\
\bottomrule
\end{tabulary}
\end{table}

科技词汇尤其依赖词义引申造出新义：web(蜘蛛网$\rightarrow$万维网)、cloud(云$\rightarrow$云计算)、window(窗$\rightarrow$窗口)、desktop(桌面$\rightarrow$桌面系统)、virus(病毒$\rightarrow$计算机病毒)。

\subsubsection{比喻造词:隐喻与转喻}

比喻造词通过隐喻(metaphor)与转喻(metonymy)把具体概念投射到抽象领域，是最富创造性的构词方式：

\begin{itemize}
    \item \textbf{隐喻(Metaphor)}: 用具体事物理解抽象概念。计算机领域几乎全部使用隐喻造词：鼠标(mouse)、病毒(virus)、防火墙(firewall)、垃圾邮件(spam)、窗口(window)、桌面(desktop)、云计算(cloud)、瓶颈(bottleneck)。汉语"鼠标"正是"鼠+标"的隐喻造词。
    \item \textbf{转喻(Metonymy)}: 用相关事物代替目标事物。白宫(代美国政府)、华盛顿(代美国政府)、华尔街(代金融界)、皇冠(代王权)、"红领巾"(代少先队员)。
    \item \textbf{提喻(Synecdoche)}: 用部分代整体(或反之)。"人手"(hand 代工人)、"饭碗"(代工作)。
\end{itemize}

认知语言学认为，人类抽象思维本身就建立在隐喻之上("时间是金钱""争论是战争")，语言中的比喻造词只是这一认知机制的表层体现。

\subsection{借词:外来词的吸收}

语言接触导致的词汇借入(borrowing)是最重要的"外部"造词来源。外来词进入汉语时，要接受汉语语音、语法和构词规则的改造，形成多种吸收方式：

\begin{table}[H]
\centering
\caption{外来词的吸收方式}
\begin{tabulary}{\textwidth}{@{}CCC@{}}
\toprule
\textbf{方式} & \textbf{机制} & \textbf{例子} \\
\midrule
音译 & 借音不借形 & 沙发(sofa)、咖啡(coffee) \\
意译 & 译义造新词 & 电脑(computer)、足球(football) \\
仿译 & 逐语素对应翻译 & 热狗(hot dog)、黑板(blackboard) \\
音义兼译 & 音义兼顾 & 可口可乐(Coca-Cola)、黑客(hacker) \\
借形 & 借字形不借音 & 电话、经济(自日语) \\
语义借用 & 旧词借新义项 & 病毒($\rightarrow$计算机病毒) \\
\bottomrule
\end{tabulary}
\end{table}

\subsubsection{音译(Transliteration)}

音译(Transliteration)按源语读音用汉字转写，汉字只表音、不表义：
\begin{itemize}
    \item \textbf{纯音译}: 沙发(sofa)、咖啡(coffee)、巧克力(chocolate)、的士(taxi)、比基尼(bikini)
    \item \textbf{音译+类名}: 在音译后加汉语类名帮助理解，如"卡车"(car+车)、"啤酒"(beer+酒)、"芭蕾舞"(ballet+舞)、"艾滋病"(AIDS+病)
    \item \textbf{半音半意}: 前音译后意译，如"冰淇淋"(ice-cream)、"呼啦圈"(hula-hoop)、"因特网"(internet)
\end{itemize}

音译词可能因方言读音不同而有异体，如 chocolate 普通话作"巧克力"、粤语作"朱古力"；人名、地名则按国家规定的音译规则处理(如 Clinton$\rightarrow$克林顿)。

\subsubsection{意译(Semantic Translation)}

意译用本族语材料把外来概念\emph{整体}翻译成新词，不保留源语读音：
\begin{itemize}
    \item 电脑(computer)、软件(software)、硬件(hardware)、数据库(database)
    \item 互联网(Internet)、搜索引擎(search engine)、足球(football)、口香糖(chewing gum)
\end{itemize}

意译词在汉语外来词中占相当比重：汉语双音节为主、表意文字的特点，使音译词(如"德律风")往往被意译词(如"电话")淘汰。意译与仿译的区别在于：意译只翻译整体概念，仿译还要保留源语的构词结构(见下)。

\subsubsection{仿译(Calque)}

仿译(Calque，又称借译、仿造)把源语词的\emph{每个语素}逐一翻译，保留源语的内部结构：
\begin{itemize}
    \item 热狗(hot dog)、黑板(blackboard)、马力(horsepower)、超人(superman)
    \item 冷战(cold war)、蓝图(blueprint)、瓶颈(bottleneck)、超级市场(supermarket)
    \item 千年虫(millennium bug)、防火墙(firewall)
\end{itemize}

仿译与意译的界限并不总是清晰："足球"(football)既可以说是意译(整体翻译概念)，也符合仿译的定义(逐语素对应 foot=足、ball=球)。一般以"是否逐语素保留源语结构"来区分——热狗是仿译(因为 hot dog 作为食物名与字面义无关，只能是逐成分仿造)，电脑则是意译(computer 没有可拆解的语素对应)。

\subsubsection{音义兼译(Phono-Semantic Matching)}

音义兼译所选汉字既接近源语读音，又能在意义上产生联想，是"形神兼备"的理想译法：
\begin{itemize}
    \item 可口可乐(Coca-Cola)、黑客(hacker)、维他命(vitamin)、香波(shampoo)
    \item 幽默(humour)、逻辑(logic)、俱乐部(club)、奔驰(Benz)、宝马(BMW)
    \item 托福(TOEFL)、雅思(IELTS)
\end{itemize}

"可口可乐"既是 Coca-Cola 的音译，又暗示"可口、快乐"；"奔驰"既近 Benz 读音又暗示汽车飞驰；"宝马"既近 BMW 的读音又令人联想到名马。这类译名音义兼得，是语言学界公认的最佳译法。

\subsubsection{借形(Borrowing of Written Form)}

借形直接借用源语的\emph{书写形式}，读音按本族语，是汉字文化圈特有的借词方式。最典型的是汉语借自日语的"和制汉语词"——日本在明治维新时期用汉字创造或重新组合的大量译词，随后回流中国：
\begin{itemize}
    \item 汉语借自日语的"和制汉语": 电话、经济、社会、干部、哲学、法律、美术、物理、化学、文化、主义、政党、义务
    \item 英语借自汉语(借音): tofu(豆腐)、kung fu(功夫)、dim sum(点心)、typhoon(台风)、tea(茶)
\end{itemize}

借形与借音不同：英语借 tofu 属于音译(借读音)，汉语借"经济"只借字形、读音完全是汉语的。由于日语词义也用汉字承载，汉语借形词几乎"无缝"融入词汇系统，成为现代汉语的基本词汇。

\subsubsection{语义借用(Semantic Loan)}

语义借用不借字形、不借读音，只把源语词的一个\emph{新义项}加在本族已有的词上：
\begin{itemize}
    \item 病毒：本义"生物病毒"，借英语 virus 的引申义后增加"计算机病毒"
    \item 明星：本义"明亮的星"，借英语 star 的引申义后增加"演艺明星"
    \item 天使：本义"神的使者"，借英语 angel 的引申义后增加"善良的人"
    \item 人口爆炸、信息爆炸：仿译 population explosion 的同时，"爆炸"获得"急剧增加"的新义
\end{itemize}

语义借用与仿译的区别：仿译产生新词(热狗是新的组合)，语义借用不产生新词、只是旧词获得新义。判断一个"新义"是自身引申还是语义借用，要看语言接触的历史。

\begin{example}[同一外来概念可有不同吸收方式]
\begin{align}
    \text{Coca-Cola} &\rightarrow \text{可口可乐(音义兼译)} \qquad \text{laser} \rightarrow \text{莱塞(音译)} / \text{激光(意译)} \\
    \text{telephone} &\rightarrow \text{德律风(音译)} / \text{电话(借形自日语)} \\
    \text{taxi} &\rightarrow \text{的士(音译)} / \text{出租车(意译)}
\end{align}
同一外来词往往音译与意译并存，最后由语言社区的选择淘汰其一部。
\end{example}

\subsection{汉语的造字法:六书}

汉语的构词基础是汉字的构造。东汉许慎在《说文解字》中提出"六书"理论，系统归纳了汉字的构造与使用方式。其中象形、指事、会意、形声是\emph{造字法}，转注、假借是\emph{用字法}：

\begin{table}[H]
\centering
\caption{六书:汉字的构造方式}
\begin{tabulary}{\textwidth}{@{}CCCC@{}}
\toprule
\textbf{名称} & \textbf{性质} & \textbf{说明} & \textbf{字例} \\
\midrule
象形 & 造字法 & 直接描画事物形状 & 日月山水 \\
指事 & 造字法 & 抽象符号/在象形上加点标示 & 上下刃本 \\
会意 & 造字法 & 两个以上部件组合出新义 & 武信休森 \\
形声 & 造字法 & 形符表义+声符表音 & 江河清湖 \\
转注 & 用字法 & 同义字互相注释 & 考老 \\
假借 & 用字法 & 借同音字记录新义 & 令长其 \\
\bottomrule
\end{tabulary}
\end{table}

\begin{itemize}
    \item \textbf{象形(Pictographic)}: 描摹客观事物的形状。"日"像太阳、"月"像月牙、"水"像流水、"山"像山峰。象形是最古老的造字法，但数量很少。
    \item \textbf{指事(Indicative)}: 用抽象符号标示位置或状态。"上""下"用一长横加一短横标示方位，"刃"在"刀"的象形上加一点，指出刀刃所在。
    \item \textbf{会意(Associative compound)}: 用两个或以上部件会合出新义。"休"=人+木(人倚树休息)、"明"=日+月(明亮)、"森"=三木(树木众多)。
    \item \textbf{形声(Phonetic compound)}: 由表义的形符和表音的声符构成，是汉字中最能产的造字法，占现代汉字的80\%以上。"江"从水工声、"河"从水可声、"清"从水青声。近代化学元素"氦""氧""氢"等都是新造的形声字。部分声符还兼表意义(如"悌"的声符"弟"兼提示"弟弟敬爱兄长"义)，称"形声兼会意"。
    \item \textbf{转注(Mutual explanatory)}: 同义字互相注释。"考""老"同义，许慎以"考"训"老"、以"老"训"考"。学界对转注历来有"互训说""部首说""语音说"等不同解释，尚无定论，一般归为用字法。
    \item \textbf{假借(Phonetic loan)}: 借用已有的同音字记录新词义。"令"本义命令，借作县令之"令"；"长"本义长久，借作官长之"长"；"其"本义簸箕，借作代词"其"。假借也是用字法，不造新字。
\end{itemize}

\subsection{构词方式的跨语言对比}

不同语言在构词方式的侧重上有明显差异，这本身也是跨语言比较的重要维度：

\begin{table}[H]
\centering
\caption{不同语言构词方式的侧重}
\begin{tabulary}{\textwidth}{@{}CCCC@{}}
\toprule
\textbf{方式} & \textbf{汉语} & \textbf{英语} & \textbf{日语} \\
\midrule
复合 & 极能产(火车) & 极能产(blackboard) & 极能产(山道) \\
派生 & 较少(词缀有限) & 极能产(-ness, re-) & 极能产(动词变形丰富) \\
混成 & 较少 & 能产(brunch) & 罕见 \\
截短/缩写 & 较少(北外) & 能产(lab, FBI) & 常见(スマホ) \\
借词 & 音译+意译+借形 & 直接借用(sofa) & 外来语音译+借形汉字词 \\
六书造字 & 独有 & — & — \\
\bottomrule
\end{tabulary}
\end{table}

这一节说明：\emph{词的创造有规律可循}。复合、派生等"组合型"方式在汉语、英语、日语中都极其能产；缩短型方式(截短、缩写、混成)在英语中尤为活跃；而借词方式则与书写系统和文化接触密切相关——汉语的六书造字法与丰富的音译/意译/借形方式，正是汉字文化圈与外部语言长期接触的产物。

\subsection{形态音位学}

语素组合时常常发生语音变化，研究这种变化的学科叫\emph{形态音位学}(Morphophonology)。最经典的例子是英语复数后缀的三种读音。

\subsubsection{英语复数规则}

\begin{equation}
    \text{复数后缀} = 
    \begin{cases}
        /s/ & \text{清辅音后 (cats)} \\
        /z/ & \text{浊辅音/元音后 (dogs, bees)} \\
        /\ipa{ɪ}z/ & \text{咝音后 (buses, churches)}
    \end{cases}
\end{equation}

形式化表示(这是音系规则格式的典型应用)：
\begin{equation}
    \text{Plural} \rightarrow 
    \begin{cases}
        [s] / [\text{-voice, -sibilant}] \_ \\
        [z] / [\text{+voice}] \_ \text{ 或 } [\text{+sonorant}] \_ \\
        [\ipa{ɪ}z] / [\text{+sibilant}] \_
    \end{cases}
\end{equation}

这个规则不是英语独有的"拼写习惯"，而是语音学上的普遍倾向——清音与清音相邻、浊音与浊音相邻(语音同化)。在土耳其语、德语、俄语的复数/格后缀中可以看到完全相同的机制。

\subsubsection{语素变体(Allomorph)}

复数语素 \{-s\} 在 cats、dogs、buses 中有三个变体 /s/、/z/、/\ipa{ɪ}z/，它们是同一个语素的\emph{语素变体}。异干法(go/went)是语素变体的极端情况——词根形式完全互补。

\subsection{形态分析的形式化}

\subsubsection{有限状态形态学}

形态分析可以建模为\emph{有限状态转换器}(FST, Finite-State Transducer)：一个把词的表层拼写映射为"词干+语法标注"序列的有限自动机。
\begin{equation}
    \text{FST: } T \subseteq \Sigma^* \times \Sigma^*
\end{equation}

\begin{example}
\begin{align*}
    \text{输入: } & \text{cats} \\
    \text{输出: } & \text{cat +N +Pl}
\end{align*}
FST把表层形式 \emph{cats} 转写为词汇层分析 \emph{cat + N + Pl}。形态学自动机是许多NLP系统(拼写检查、词形还原、机器翻译形态分析)的基础。
\end{example}

\subsubsection{双层面模型}

著名的双层面形态学(Koskenniemi, 1983)把形态变化表示为两个层面的对应：
\begin{equation}
    \text{词表层形式} \xleftrightarrow{\text{FST}} \text{词法分析}
\end{equation}

规则示例：
\begin{equation}
    \begin{aligned}
    &\text{Lexical: } cat +N +Pl \\
    &\xrightarrow{\text{拼写规则}} \\
    &\text{Surface: } cats
    \end{aligned}
\end{equation}

\begin{example}[y→i拼写规则]
\begin{equation}
    \text{carry + Pl} \xrightarrow{y \rightarrow ies} \text{carries}
\end{equation}
双层面模型中，词汇层面是 \emph{carry+Pl}，表层是 \emph{carries}，FST同时处理"插入e"和"y变i"两个操作。
\end{example}

\subsection{构词规则的形式化}

\subsubsection{语素顺序约束}

英语词缀有严格的线性顺序，这是形态学中的普遍约束：
\begin{equation}
    \text{词根} \rightarrow \text{派生前缀} \rightarrow \text{词根} \rightarrow \text{派生后缀} \rightarrow \text{屈折后缀}
\end{equation}

\begin{example}
\begin{equation}
    \text{un-happi-ness-es}
\end{equation}
\begin{itemize}
    \item un-: 派生前缀
    \item happi: 词根
    \item -ness: 派生后缀
    \item -es: 屈折后缀
\end{itemize}
这个顺序不是偶然的：屈折后缀永远在最外层(英语中动词屈折不能出现在派生词缀之内)。汉语中同样存在顺序约束："桌子们"合法而"*桌子"的复数须加"们"且不能提前。
\end{example}

\subsubsection{形态生成的图论表示}

词的形态结构可以看作一棵树(形态树)：词根是根节点，词缀是叶子：
\begin{equation}
    \text{un-happiness: } \underbrace{un}_{\text{前缀}} + \underbrace{happi}_{\text{词根}} + \underbrace{ness}_{\text{后缀}}
\end{equation}

\subsection{形态类型学的形式刻画}

上一章的孤立语/黏着语/屈折语/多式综合语四分法，可以在形态学的框架下精确刻画。定义语素密度(词素比率)为：
\begin{equation}
    \text{形态综合度} = \frac{\#\text{语素}}{\#\text{词}} \quad (\text{对一段语料取平均})
\end{equation}

\begin{itemize}
    \item 孤立语(汉语): 比率 $\approx 1$
    \item 黏着语(土耳其语): 比率高且词缀边界清晰(一一对应)
    \item 屈折语(俄语): 比率较高且词缀融合(一个词缀多重意义)
    \item 多式综合语(因纽特语): 比率极高(一个词内整合主语宾语)
\end{itemize}

这个量化指标是跨语言比较的直接工具：给定两段平行文本，计算各自的语素密度，就能客观比较两种语言在形态上的"综合性"差异。

\subsection{本章小结}

形态学研究词内的结构世界：语素是最小的意义单位，屈折与派生是两大构词机制，语素变体和语音规则把深层形式映射到表层拼写。形态学既能用有限状态转换器精确计算，又是语言类型学(孤立/黏着/屈折/复综)的核心坐标。不同语言在形态上的巨大差异——从汉语的几乎没有屈折，到俄语的复杂格系统，到日语的黏着堆叠——正是理解"语言之间差异"的最佳入口，我们将在最后一章系统比较。下一章上升到句子的组织方式：句法。

\newpage
\section{句法}

\subsection{句法概述}

句法(Syntax)研究词如何组合成短语和句子的规则。句法不仅告诉我们\emph{什么合法}(合法句子)，也告诉我们\emph{为什么合法}——合法性与句子结构的层次性、递归性紧密相关。

\begin{example}[句法的层级性]
\begin{equation}
    \text{The cat} \quad \text{在句中作为整体出现} \quad \rightarrow \quad \text{The cat sat on the mat}
\end{equation}
"the cat"这个名词短语(NP)可以出现在"\_\_\_ sleeps""I saw \_\_\_""\_\_\_ is big"等位置，说明句法是成分组成的层级结构，而不是词的线性序列。
\end{example}

句法的核心问题：
\begin{itemize}
    \item 短语结构如何组织(成分分析)
    \item 词类(范畴)如何决定组合方式
    \item 语序为何因语言而异(跨语言比较的关键)
    \item 有限的规则如何生成无限的句子(递归性)
\end{itemize}

\subsection{短语结构语法}

\subsubsection{基本短语规则}

短语结构语法用上下文无关文法式的重写规则描述句子结构：
\begin{equation}
    \begin{aligned}
    S &\rightarrow NP \ VP \\
    NP &\rightarrow Det \ (Adj) \ N \ (PP) \\
    VP &\rightarrow V \ (NP) \ (PP) \ (Adv) \\
    PP &\rightarrow P \ NP \\
    AdjP &\rightarrow (Adv) \ Adj \\
    AdvP &\rightarrow (Adv) \ Adv
    \end{aligned}
\end{equation}

这里 $S$(句子)、$NP$(名词短语)、$VP$(动词短语)、$PP$(介词短语)是非终结符，规则递归地组合出无限多的句子。注意$VP \rightarrow V\ (NP)$ 允许多种结构的递归嵌套——这正是有限规则生成无限句子的机制。

\subsubsection{句法树示例}

"The cat sat on the mat":

\begin{forest}
for tree={s sep=8mm, l sep=8mm}
[S
    [NP
        [Det [The]]
        [N [cat]]
    ]
    [VP
        [V [sat]]
        [PP
            [P [on]]
            [NP
                [Det [the]]
                [N [mat]]
            ]
        ]
    ]
]
\end{forest}

句法树的每个子树对应一个\emph{成分}(constituent)，成分性是短语结构语法的基本概念：一个句子可以"被测试"出成分(前置、替代、省略测试)。

\subsubsection{短语结构语法的类型学参数}

同一套重写规则在不同语言中只是\emph{语序不同}。比较英语与日语：
\begin{align}
    \text{英语: } & S \rightarrow NP \ VP \quad (\text{中心词在前: 动-宾}) \\
    \text{日语: } & S \rightarrow NP \ VP \quad (\text{中心词在后: 宾-动}) \\
    \text{日语: } & NP \rightarrow NP \ P \quad (\text{后置词})
\end{align}
短语结构"骨架"惊人地一致，差异集中在中心词的方向——这一观察直接导向X-bar理论中的头参数。

\subsection{X-bar理论}

X-bar理论(Chomsky, 1970)提供统一的短语结构模板，消除了为每种短语分别写规则的必要。每个短语XP都有一个中心词X，模板如下：
\begin{equation}
    \begin{aligned}
    XP &\rightarrow (Spec) \ X' \\
    X' &\rightarrow X \ (Complement) \\
    X' &\rightarrow (Adjunct) \ X'
    \end{aligned}
\end{equation}

\begin{forest}
for tree={s sep=8mm, l sep=8mm}
[XP
    [Spec]
    [X$'$
        [X$'$
            [Adjunct]
            [X$'$
                [X]
                [Complement]
            ]
        ]
    ]
]
\end{forest}

X-bar统一刻画了名词短语NP、动词短语VP、介词短语PP、形容词短语AP等一切短语：如NP的中心词是N，VP的中心词是V。通过递归的X'层，X-bar允许任意多个修饰语(Adjunct)附着，同时保持每个短语只有一个中心词和一个补足语。

\subsubsection{头参数与跨语言差异}

\begin{example}[头参数]
\begin{align}
    &\text{英语(中心词在前):} \quad NP = P + NP \quad (\text{前置词: in the house}) \\
    &\text{日语(中心词在后):} \quad NP = NP + P \quad (\text{后置词: 家 の 中}) \\
    &\text{英语(动-宾):} \quad VP = V + NP \quad (\text{eat an apple}) \\
    &\text{日语(宾-动):} \quad VP = NP + V \quad (\text{りんご を 食べる})
\end{align}
头参数(Head parameter)决定了中心词在其补足语之前还是之后。这一参数贯穿整个短语结构，解释了英语(中心词前)与日语(中心词后)在语序上的系统性差异。我们将在最后一章详细展开。
\end{example}

\subsection{依存语法}

依存语法(Dependency Grammar)不关心成分层级，而是关注\emph{词与词之间不对称的二元关系}。每个词有一个中心词(head)，全句的中心词是动词：

\begin{equation}
    \text{依存树} = (W, E)
\end{equation}
其中 $W$ 是词集合，$E \subseteq W \times R \times W$ 是有标记的依存边。

\subsubsection{依存关系类型}

\begin{table}[H]
\centering
\caption{常见依存关系(以英语为例)}
\begin{tabulary}{\textwidth}{@{}CC@{}}
\toprule
\textbf{关系} & \textbf{示例} \\
\midrule
nsubj (名词主语) & \textbf{cat} sleeps \\
dobj (直接宾语) & eat \textbf{apple} \\
amod (形容词修饰) & \textbf{red} car \\
det (限定词) & \textbf{the} cat \\
prep (介词修饰) & sit \textbf{on} mat \\
pobj (介词宾语) & on \textbf{mat} \\
advmod (副词修饰) & \textbf{very} good \\
aux (助动词) & \textbf{has} eaten \\
\bottomrule
\end{tabulary}
\end{table}

\subsubsection{依存语法的跨语言优势}

依存语法对语序高度"宽容"：不论动词在宾语前(SVO)还是后(SOV)，\emph{动词-宾语}的依存关系都成立。这使依存表示成为跨语言句法分析(Universal Dependencies项目)的标准格式——同一棵"语义骨架"在不同语言中有不同的表层语序，依存树完美地分离了"关系"与"语序"。这与短语结构语法形成互补：短语结构刻画成分层级，依存结构刻画中心词-从属关系。

\subsection{转换生成语法}

\subsubsection{深层结构与表层结构}

早期生成语法(1957-1980年代)区分深层结构与表层结构：
\begin{equation}
    \text{深层结构(DS)} \xrightarrow{\text{转换规则}} \text{表层结构(SS)}
\end{equation}

深层结构承载论元结构与语义解释，表层结构是实际说出的形式。转换规则(移动、插入、删除)把深层结构映射到表层。这一机制使生成语法可以优雅地处理"同一个意思的多种句子形式"(主动/被动、陈述/疑问)之间的关系。

\subsubsection{Wh-移动}

\begin{align}
    \text{DS: } & \text{You saw } \underline{\text{what}} \\
    \text{SS: } & \text{What}_i \text{ did you see } t_i
\end{align}

\begin{equation}
    \text{Move-}\alpha: \text{将成分 } \alpha \text{ 移至指定位置，留下语迹 } t
\end{equation}

疑问词 \emph{what} 从宾语位置移动到句首，原位置留下语迹 $t_i$(下标的 $i$ 表示同标索引，标明移动轨迹)。语迹理论解释了疑问句为何不能从某些位置提取(如孤岛约束)。

\subsubsection{孤岛约束}

某些结构是"移动禁区"：成分不能从中提取。
\begin{equation}
    \text{*What did you ask whether John bought \_ ?} \quad (\text{whether 从句是孤岛})
\end{equation}
这类约束在不同语言中表现出惊人的一致性，说明它们源于语法的普遍原则(Subjacency/条件约束)，而非个别语言的规则——这是普遍语法的重要证据。

\subsection{最简方案}

\subsubsection{合并(Merge)}

最简方案(Chomsky, 1995-)把语法简化到极致：句法操作只剩下\emph{合并}(Merge)。Merge把两个句法对象组合成一个新对象：
\begin{equation}
    \text{Merge}(\alpha, \beta) = \{\alpha, \{\alpha, \beta\}\}
\end{equation}

合并是\emph{递归}的，因而能生成无限结构：
\begin{equation}
    \begin{aligned}
    &\text{Merge}(the, cat) = \{the, \{the, cat\}\} = NP \\
    &\text{Merge}(NP, saw) = \{NP, \{NP, saw\}\} = VP \\
    &\text{Merge}(VP, NP) = \{VP, \{VP, NP\}\} = S
    \end{aligned}
\end{equation}

\subsubsection{移动是合并的副产品}

最简方案中，移动(如Wh-移动)只是"内部合并"(再合并已经合并过的成分)。语言的递归性和"离散无限性"(离散的、无限的能力)由此获得最简约的解释：
\begin{equation}
    \text{语言的生成能力} = \text{合并的递归应用}
\end{equation}

\subsection{范畴语法}

\subsubsection{组合范畴语法(CCG)}

范畴语法把每个词与一个\emph{范畴}关联，组合规则是纯粹的"函数应用"。范畴用斜线标注参数方向：
\begin{equation}
    \begin{aligned}
    \text{cat} &:= N \\
    \text{the} &:= NP/N \\
    \text{sleeps} &:= S\backslash NP \\
    \text{loves} &:= (S\backslash NP)/NP
    \end{aligned}
\end{equation}

\subsubsection{组合规则}

\begin{equation}
    \begin{aligned}
    X/Y \quad Y &\Rightarrow X \quad \text{(前向应用)} \\
    Y \quad X\backslash Y &\Rightarrow X \quad \text{(后向应用)}
    \end{aligned}
\end{equation}

\begin{example}[CCG推导]
\begin{equation}
    \frac{\text{the} : NP/N \quad \text{cat}: N}{NP} \quad (\text{前向应用})
\end{equation}
"the cat"的范畴 $NP/N$ 与 $N$ 结合，前向应用得到 $NP$。CCG的范畴系统与语义学中的类型论(最后一章的Lambda演算)完全同构：每个范畴对应一个语义类型，句法推导直接对应语义组合。
\end{example}

CCG的一个重要优点是它的组合规则能产生\emph{松散成分}(non-constituent)协调，使它在机器翻译和跨语言句法分析中特别有用。

\subsection{树邻接文法(TAG)}

TAG使用两棵树作为基本对象：
\begin{equation}
    G_{TAG} = (\Sigma, V, I, A, S)
\end{equation}
其中 $I$ 是初始树集合，$A$ 是辅助树集合。

操作：
\begin{itemize}
    \item \textbf{替换(Substitution)}: 将一棵树插入叶节点
    \item \textbf{附加(Adjoining)}: 将辅助树插入内部节点(可以描述长距离依赖和修饰语的递归嵌套)
\end{itemize}

TAG的表达能力恰好处于上下文无关文法与上下文有关文法之间(\emph{温和上下文有关})，能处理自然语言中的跨串依赖，同时保持多项式可解析性。这是TAG在计算语言学中占据重要地位的原因。

\subsection{概率句法分析}

\subsubsection{概率上下文无关文法(PCFG)}

把概率赋给每条产生式，就得到PCFG。它允许在歧义句法树之间选择概率最高的一棵：
\begin{equation}
    P(T|S) = \prod_{r \in T} P(r)
\end{equation}

其中 $T$ 是句法树，$S$ 是句子，$r$ 是产生式规则。PCFG可以从语料自动学习(最大似然估计)，是统计句法分析的基础。

\subsubsection{CKY算法}

CKY算法是PCFG句法分析的标准动态规划算法(要求文法为CNF)：
\begin{equation}
    \pi[i,j,X] = \max_{Y,Z,k} P(X \rightarrow YZ) \cdot \pi[i,k,Y] \cdot \pi[k+1,j,Z]
\end{equation}

\begin{equation}
    \pi[i,i,X] = P(X \rightarrow w_i)
\end{equation}

CKY自底向上填表：先算单个词的可能范畴，再逐步合并相邻区间，最终检查 $\pi[1,n,S]$ 是否为最大概率。对长度为 $n$ 的句子，复杂度为 $O(n^3)$。CKY在词性标注、机器翻译、语音识别的句法约束中被广泛使用。

\subsection{跨语言的句法参数}

句法理论积累的洞见可以概括为一组跨语言比较的参数(详见最后一章)：
\begin{itemize}
    \item \textbf{头参数}: 中心词在前(英语) vs. 中心词在后(日语、土耳其语)
    \item \textbf{主语优先 vs. 话题优先}: 英语、日语是主语优先，汉语、朝鲜语是话题优先
    \item \textbf{Pro-省略(Pro-drop)}: 意大利语、西班牙语、日语可省略主语，英语、法语、德语不可
    \item \textbf{疑问词移动 vs. 原位}: 英语疑问词移动，汉语、日语疑问词留在原位
    \item \textbf{格标记 vs. 语序}: 俄语靠格标记，汉语靠语序
\end{itemize}

\begin{example}[疑问词原位 vs 移动]
\begin{align}
    \text{英语: } & \text{What did you see \_ ?} \quad (\text{移动}) \\
    \text{汉语: } & \text{你看到了什么?} \quad (\text{原位}) \\
    \text{日语: } & \text{あなたは何を見たか?} \quad (\text{原位})
\end{align}
同一命题"你看见了什么"在三语言中的句法实现截然不同：英语把疑问词提到句首，汉语和日语留在宾语位置。这是"参数差异"最直观的例子之一。
\end{example}

\subsection{本章小结}

句法是语言学的核心：短语结构语法刻画成分层级，X-bar理论统一短语模板并引出头参数，依存语法刻画中心词关系且天然跨语言友好，转换语法与最简方案追问"句法操作的普遍形式"，范畴语法与TAG提供计算上可实现的句法形式体系，PCFG与CKY让句法分析可以自动化。句法的深层启示是：人类语言共享一套普遍原则(递归合并、中心词-补足语结构)，而语言之间的差异源于少数参数的取值。这正是最后一章"不同语言之间的差异"的理论基础。下一章转向意义：语义学。

\newpage
\section{语义学}

\subsection{语义学概述}

语义学(Semantics)研究语言单位——词、短语、句子——的意义。形式语义学的基本信念是：\emph{意义可以用数学对象刻画}，句子的意义可以用真值条件(在什么条件下为真)来把握。

\begin{equation}
    \text{知道一个句子的意义} = \text{知道它在什么条件下为真}
\end{equation}

语义学与语用学(Pragmatics)的分界：语义学研究"句子说了什么"(不依赖语境的字面意义)，语用学研究"说话人想表达什么"(依赖语境的含义)。例如"你能把盐递过来吗"字面意义是疑问，语用意义是请求。

\subsection{词义关系}

\subsubsection{基本关系}

\begin{itemize}
    \item \textbf{同义(Synonymy)}: 意义相同或几乎相同，big $\approx$ large
    \item \textbf{反义(Antonymy)}: 意义相反。分互补反义(hot/cold 非此即彼式不严格)与程度反义，还有方向反义(buy/sell)
    \item \textbf{上下义(Hyponymy)}: 从属关系，dog $\subset$ animal，dog是animal的下义词
    \item \textbf{整体部分(Meronymy)}: 部分-整体关系，wheel $\subset$ car
    \item \textbf{多义(Polysemy)}: 一词多义且意义相关，bank(银行/河岸)
    \item \textbf{同音(Homophony)}: 发音相同意义无关，flower/flour
\end{itemize}

这些关系可以形式化为集合之间的包含、相交、矛盾关系：
\begin{equation}
    \llbracket \text{dog} \rrbracket \subseteq \llbracket \text{animal} \rrbracket \quad (\text{上下义=集合包含})
\end{equation}

\subsection{形式语义学}

\subsubsection{模型论语义}

\begin{definition}[模型]
模型 $M = \langle D, I \rangle$
\begin{itemize}
    \item $D$: 论域(Domain of discourse)——语境中的实体集合
    \item $I$: 解释函数(Interpretation function)——把专名映射到实体、把谓词映射到性质(集合)
\end{itemize}
\end{definition}

\begin{example}
模型 $M$: $D = \{ \text{猫1, 猫2, 狗1} \}$，$I(\text{cat}) = \{ \text{猫1, 猫2} \}$，$I(\text{Fluffy}) = \text{猫1}$。于是"Fluffy is a cat"在$M$中为真，因为 $I(\text{Fluffy}) \in I(\text{cat})$。
\end{example}

\subsubsection{真值条件语义}

句子意义的真值条件定义：
\begin{equation}
    \llbracket \phi \rrbracket^M = 
    \begin{cases}
        1 & \text{若 } \phi \text{ 在模型 } M \text{ 中为真} \\
        0 & \text{若 } \phi \text{ 在模型 } M \text{ 中为假}
    \end{cases}
\end{equation}

\begin{example}
\begin{align*}
    \llbracket \text{Snow is white} \rrbracket &= 1 \quad \text{当且仅当雪是白的} \\
    \llbracket \text{Cats are dogs} \rrbracket &= 0
\end{align*}
\end{example}

\subsubsection{谓词逻辑翻译}

句子翻译为一阶逻辑公式，谓词的语义即论域的子集：
\begin{equation}
    \llbracket \text{cat}(x) \rrbracket^M = 1 \iff I(x) \in I(\text{cat}) \subseteq D
\end{equation}

\begin{example}
\begin{equation}
    \text{"Every dog barks"} \Longrightarrow \forall x[\text{dog}(x) \rightarrow \text{bark}(x)]
\end{equation}
\end{example}

\subsection{Lambda演算与组合语义}

\subsubsection{Lambda演算基础}

\begin{equation}
    \begin{aligned}
    &\text{项: } M, N ::= x \mid \lambda x.M \mid (M \ N) \\
    &\beta\text{-归约: } (\lambda x.M)(N) \rightarrow M[x := N]
    \end{aligned}
\end{equation}

$\beta$-归约是计算规则：把函数应用到参数，用参数替换函数体内的绑定变量。这为语义组合提供了机械的计算手段。

\subsubsection{词项语义}

\begin{equation}
    \begin{aligned}
    \llbracket \text{cat} \rrbracket &= \lambda x.\text{cat}(x) \\
    \llbracket \text{sleep} \rrbracket &= \lambda x.\text{sleep}(x) \\
    \llbracket \text{love} \rrbracket &= \lambda y.\lambda x.\text{love}(x,y)
    \end{aligned}
\end{equation}

注意及物动词 love 是"先取宾语再取主语"的柯里化函数。语义类型与句法范畴严格对应：$N$ 对应个体($e$)，$S$ 对应真值($t$)，$NP/N$ 对应函数类型等——这就是\emph{类型-句法同构}(对应CCG的范畴系统)。

\subsubsection{组合示例}

"The cat sleeps":
\begin{equation}
    \begin{aligned}
    \llbracket \text{the} \rrbracket &= \lambda P.\iota x.P(x) \quad (\iota \text{ 是唯一性算子}) \\
    \llbracket \text{cat} \rrbracket &= \lambda x.\text{cat}(x) \\
    \llbracket \text{the cat} \rrbracket &= \iota x.\text{cat}(x) \\
    \llbracket \text{sleeps} \rrbracket &= \lambda x.\text{sleep}(x) \\
    \llbracket \text{The cat sleeps} \rrbracket &= \text{sleep}(\iota x.\text{cat}(x))
    \end{aligned}
\end{equation}

整个推导完全由句法结构(树)驱动：每个成分的意义是子成分意义按组合规则的机械应用。这是弗雷格组合原则(Compositionality)的形式实现：\emph{整体意义是部分意义+组合方式的函数}。

\subsection{量词语义}

\subsubsection{广义量词}

普通名词短语如 "every dog" 不能简单地翻译为个体，需要广义量词理论。量词是"名词谓词与动词短语谓词之间关系的函数"：
\begin{equation}
    \begin{aligned}
    \llbracket \text{every} \rrbracket &= \lambda P.\lambda Q.\forall x[P(x) \rightarrow Q(x)] \\
    \llbracket \text{some} \rrbracket &= \lambda P.\lambda Q.\exists x[P(x) \land Q(x)] \\
    \llbracket \text{no} \rrbracket &= \lambda P.\lambda Q.\neg\exists x[P(x) \land Q(x)] \\
    \llbracket \text{most} \rrbracket &= \lambda P.\lambda Q.|P \cap Q| > |P \setminus Q|
    \end{aligned}
\end{equation}

广义量词是"集合之间的关系"，因此可以统一地表示 "most dogs" 这种一阶逻辑难以直接表达的量词。

\subsubsection{量词作用域}

量词之间的相对顺序产生不同的解读，导致歧义：
\begin{equation}
    \begin{aligned}
    &\text{Every student read a book} \\
    &\text{解读1 (每一>存在): } \forall x[\text{student}(x) \rightarrow \exists y[\text{book}(y) \land \text{read}(x,y)]] \\
    &\text{解读2 (存在>每一): } \exists y[\text{book}(y) \land \forall x[\text{student}(x) \rightarrow \text{read}(x,y)]]
    \end{aligned}
\end{equation}

解读1:"每个学生读了一本书(可能是不同的书)"；解读2:"有一本书，每个学生都读了它(同一本书)"。形式语义学用作用域(scope)解释这种歧义，量词作作用域越广(越靠左的算子)，辖域越大。

\subsection{内涵语义}

\subsubsection{可能世界语义}

外延语义无法处理"相信""可能"等内涵语境。可能世界语义学(Kripke, 1963)引入可能世界的概念：
\begin{equation}
    \llbracket \phi \rrbracket^{M,w,g} \in \{0, 1\}
\end{equation}
其中 $w$ 是可能世界，$g$ 是赋值函数。

\subsubsection{模态算子}

\begin{equation}
    \begin{aligned}
    \llbracket \Box \phi \rrbracket^{M,w} &= 1 \text{ 当且仅当 } \forall w' \in R(w): \llbracket \phi \rrbracket^{M,w'} = 1 \\
    \llbracket \Diamond \phi \rrbracket^{M,w} &= 1 \text{ 当且仅当 } \exists w' \in R(w): \llbracket \phi \rrbracket^{M,w'} = 1
    \end{aligned}
\end{equation}

$\Box$(必然)要求 $\phi$ 在所有可达世界中为真，$\Diamond$(可能)要求 $\phi$ 在某个可达世界中为真。可达关系 $R$ 可以取不同的性质(自反、传递等)，对应不同的模态逻辑系统(T、S4、S5等)。

\subsubsection{命题态度}

\begin{equation}
    \text{"John believes that p"} = \text{John与命题 } \lambda w.\llbracket p \rrbracket^w \text{ 处于"相信"关系}
\end{equation}
内涵语境(相信、希望、知道)中的从句指涉的不是真值而是命题(从世界到真值的函数)。这解释了为什么 "John believes that the morning star is the evening star" 中两个名称不能随便互换——它们的外延相同(金星)但内涵(命题角色)不同。

\subsection{时间与体态语义}

\subsubsection{时态与时间论域}

时态把时间引入语义。用时间参数化的模型：
\begin{equation}
    \llbracket \text{walked} \rrbracket^{M,g,t} = 1 \iff \text{主体在 } t' \text{ 行走，且 } t' < t
\end{equation}

\subsubsection{体态(Aspect)}

体态刻画动作的内部时间结构：
\begin{itemize}
    \item \textbf{完成体(Perfective)}: 把事件看作完整整体，"他读了那本书"
    \item \textbf{进行体(Imperfective)}: 从内部看事件进行中，"他正在读那本书"
\end{itemize}
汉语通过助词"了/着/过"表达体貌，英语用形态(-ed/-ing)表达时体，斯拉夫语(俄语)把完成/未完成固化为动词的语法对(perfective/imperfective)。这是跨语言差异在语义层面的体现。

\subsection{词汇语义表示}

\subsubsection{框架语义}

Fillmore的框架语义学用"语义框架"描述词义：一个词激活一个包含参与者角色(槽位)的场景。
\begin{equation}
    \text{Buy: } \langle \text{Buyer}, \text{Seller}, \text{Goods}, \text{Money} \rangle
\end{equation}

\begin{example}
\begin{align}
    &\text{John bought the car from Mary for \$500} \\
    &\text{Buyer=John, Goods=the car, Seller=Mary, Money=\$500}
\end{align}
同一场景可以用不同动词的视角表述：\emph{buy}(买者视角)、\emph{sell}(卖者视角)、\emph{pay}(价钱视角)。框架语义把词汇语义与句法论元结构连接起来，是FrameNet等资源的理论基础。
\end{example}

\subsubsection{论元结构(Theta角色)}

动词的语义论元映射到句法位置：
\begin{equation}
    \text{give}: \langle \text{施事(Agent)}, \text{受事(Theme)}, \text{接受者(Recipient)} \rangle
\end{equation}

\begin{equation}
    \text{John} \xrightarrow{Agent} \text{gave} \quad \text{a book} \xrightarrow{Theme} \text{to Mary} \xrightarrow{Recipient}
\end{equation}

论元结构在不同语言中体现为格系统或语序：俄语用与格(dative)标记接受者，汉语用"给"介宾结构，日语用助词"に"标记。动词的论元结构是"语义与句法的接口"。

\subsubsection{语用与语义的边界:预设与含义}

\begin{itemize}
    \item \textbf{蕴含(Entailment)}: 逻辑必然，A蕴含B指只要A真B必真
    \item \textbf{预设(Presupposition)}: 隐含前提，"他的哥哥来了"预设"他有哥哥"
    \item \textbf{会话含义(Implicature)}: 格里斯的语用推理，"有些学生来了"常暗示"不是全部"
\end{itemize}

\subsection{跨语言的语义共性}

形式语义学揭示了一个重要事实：\emph{语义结构高度跨语言一致}。无论英语、汉语还是日语，都区分施事与受事、都有量词、都区分必然与可能。语言之间的差异主要不在"意义本身"，而在"意义的表层编码方式"：
\begin{itemize}
    \item 英语用量词化的名词短语 + 语序表达"谁对谁做了什么"
    \item 汉语用语序 + 话题结构表达信息结构
    \item 日语用助词 + 形态表达论元角色
    \item 俄语用格标记表达论元角色
\end{itemize}

这个洞见是最后一章的核心主题之一：普遍语义 + 参数化编码 = 语言差异。

\subsection{本章小结}

语义学研究意义，形式语义学用模型、真值条件、Lambda演算、可能世界和广义量词把意义精确化。组合原则告诉我们意义如何由部分组合，内涵语义处理相信与模态，框架语义连接词汇与论元结构。最重要的收获是：不同语言在语义层面高度共享(论元角色、量词、时态语义是普遍的)，差异集中在编码策略。接下来我们进入语音层面，看看语言在"声音"上如何不同。

\newpage
\section{音系学与语音学}

\subsection{语音学 vs 音系学}

\begin{itemize}
    \item \textbf{语音学(Phonetics)}: 研究语音的物理与生理属性(发音、声学、感知)，处理的是\emph{连续的声音}
    \item \textbf{音系学(Phonology)}: 研究语音在特定语言系统中的功能和模式，处理的是\emph{离散的单位}(音位)
\end{itemize}

语音学回答"声音是什么"(物理现实)，音系学回答"声音在一个语言中起什么作用"(系统功能)。同一物理声音在不同语言中可能有不同地位：送气塞音 [p\textsuperscript{h}] 在英语中是 /p/ 的音位变体(不区分意义)，在印地语中却与 /p/ 对立(区分意义)。这是跨语言差异的起点——\emph{每种语言从语音的连续空间中选取自己的离散系统}。

\subsection{国际音标(IPA)}

IPA提供语音的标准化符号系统，一个符号对应一个音位/音素，避免"拼写与发音不一致"的问题(英语拼写 world 中有 o 但发音没有)。

\begin{table}[H]
\centering
\caption{辅音示例}
\begin{tabulary}{\textwidth}{@{}CCCC@{}}
\toprule
\textbf{IPA} & \textbf{描述} & \textbf{例词} & \textbf{特征} \\
\midrule
/p/ & 清双唇塞音 & pat & [-voice, +bilabial, +stop] \\
/b/ & 浊双唇塞音 & bat & [+voice, +bilabial, +stop] \\
/t/ & 清齿龈塞音 & top & [-voice, +alveolar, +stop] \\
/k/ & 清软腭塞音 & cat & [-voice, +velar, +stop] \\
/s/ & 清齿龈擦音 & sit & [-voice, +alveolar, +fricative] \\
/m/ & 双唇鼻音 & mat & [+voice, +bilabial, +nasal] \\
\bottomrule
\end{tabulary}
\end{table}

\subsubsection{元音系统}

元音按舌位(高/低、前/后)和圆唇度定位。英语有12-15个元音，日语只有5个(/a i u e o/)，这是跨语言音位库存差异的典型。汉语普通话约有6个元音音位(连同声调形成声调音位)。

\subsection{区别特征}

音位可用二元特征矩阵表示，这就是\emph{区别特征理论}(雅柯布森/乔姆斯基-哈雷SPE)。特征把语音描述为"开关"的组合：
\begin{equation}
    \text{/p/} = 
    \begin{bmatrix}
    +\text{consonantal} \\
    -\text{sonorant} \\
    -\text{voice} \\
    +\text{bilabial} \\
    -\text{continuant}
    \end{bmatrix}
\end{equation}

区别特征的价值：
\begin{itemize}
    \item 用最少的特征区分一种语言的全部音位(分类的简约性)
    \item 自然类(Natural class)：共享特征集的音位在规则中表现一致
    \item 解释规则：音系规则写在一组特征上而非单个音位上
\end{itemize}

\begin{example}
英语复数规则的三种变体(/s/ /z/ /\ipa{ɪ}z/)正是由[-voice]、[+sibilant]等特征类决定的(见形态学一章)。规则"清辅音后读清音"实际上是特征[voice]的语音同化。
\end{example}

\subsection{音系规则}

\subsubsection{SPE规则格式}

音系规则的标准格式(源自《音系学模式》SPE)：
\begin{equation}
    A \rightarrow B \ / \ C \ \_ \ D
\end{equation}
读作：A在C和D之间变为B。其中 $C \_ D$ 是触发环境。

\subsubsection{英语复数规则}

\begin{equation}
    \begin{aligned}
    &\text{规则1: } \emptyset \rightarrow [\ipa{ɪ}z] \ / \ [\text{+sibilant}] \ \_ \\
    &\text{规则2: } [z] \rightarrow [s] \ / \ [\text{-voice}] \ \_ \\
    &\text{规则3: } \emptyset \rightarrow [z] \ / \ \text{elsewhere}
    \end{aligned}
\end{equation}

应用顺序(规则有序性)：
\begin{equation}
    \begin{aligned}
    &\text{/b\ipa{ʌ}s/ + Pl} \xrightarrow{R1} \text{b\ipa{ʌ}s\ipa{ɪ}z} \quad \text{(buses)} \\
    &\text{/k\ipa{æ}t/ + Pl} \xrightarrow{R3} \text{k\ipa{æ}ts} \quad \text{(cats)} \\
    &\text{/d\ipa{ɒ}\ipa{ɡ}/ + Pl} \xrightarrow{R3} \text{d\ipa{ɒ}\ipa{ɡ}z} \quad \text{(dogs)}
    \end{aligned}
\end{equation}

注意复数语素的深层形式统一取 /z/(浊音)，规则2在清辅音后把/z/清化为/s/。这种"底层形式 + 规则推导"的思路是生成音系学的核心：\emph{底层形式}存储于词库，\emph{表层形式}是规则应用的输出。

\subsection{音位与音位变体}

\begin{definition}[音位]
音位(Phoneme)是能区分意义的最小语音单位。两个音若出现在相同的环境中且能区分意义，就是不同的音位(最小对立体测试)。
\end{definition}

\begin{equation}
    \text{/p/ vs /b/} \Rightarrow \text{pat vs bat} \quad \text{(最小对立体)}
\end{equation}

\begin{definition}[音位变体]
音位变体(Allophone)是音位的条件变体，不区分意义，其分布由语音环境决定。
\end{definition}

\begin{equation}
    \text{/p/} \rightarrow 
    \begin{cases}
        [p^h] & \text{词首 (pin)} \\
        [p] & \text{/s/后 (spin)}
    \end{cases}
\end{equation}

英语中 [p\textsuperscript{h}] 和 [p] 属于同一个音位 /p/(互补分布)，说英语的人甚至察觉不到区别。而泰语、印地语中 [p\textsuperscript{h}] 和 [p] 是不同音位。\emph{音位对立/变体分布}正是跨语言音系差异的集中体现：
\begin{itemize}
    \item 英语：送气与不送气不对立(同一音位)
    \item 印地语、泰语、汉语吴语：送气与不送气对立(不同音位)
\end{itemize}

\subsection{音节结构}

音节组织规则因语言而异。音节结构：
\begin{equation}
    \text{音节} = \text{Onset} + \text{Rhyme}
\end{equation}
\begin{equation}
    \text{Rhyme} = \text{Nucleus} + \text{Coda}
\end{equation}

\begin{forest}
for tree={s sep=8mm, l sep=8mm}
[$\sigma$
    [Onset
        [/p/]
    ]
    [Rhyme
        [Nucleus
            [/\ipa{ɪ}/]
        ]
        [Coda
            [/n/]
        ]
    ]
]
\end{forest}

\subsubsection{音节模板的跨语言差异}

\begin{itemize}
    \item 汉语普通话：音节结构非常受限，最大为(C)V(N)或(C)V(\ipa{ŋ})，没有辅音丛
    \item 日语：最大为(C)(y)V(N/Q)，N是拨音、Q是促音，禁止辅音丛
    \item 英语：允许复杂辅音丛，如 \emph{strengths} /str\ipa{e}\ipa{ŋ}\ipa{θ}s/ 有连续辅音
    \item 波兰语：允许极复杂的辅音丛，如 \emph{bezwzględy}(无情的)
\end{itemize}

这些差异决定了"听感"和"拼读"上的巨大区别：日本人学英语会在辅音丛中插入元音(eskuriputo)，中国人把英语音节套进普通话模板。音系学的目标之一就是精确刻画每种语言的音节模板。

\subsection{声调与重音}

\subsubsection{声调语言}

声调(音高模式)在声调语言中区分词义。汉语普通话四声：
\begin{equation}
    \begin{aligned}
    &\text{妈 } \ipa{mā} \ [55] \quad \text{高平} \\
    &\text{麻 } \ipa{má} \ [35] \quad \text{高升} \\
    &\text{马 } \ipa{mǎ} \ [214] \quad \text{降升} \\
    &\text{骂 } \ipa{mà} \ [51] \quad \text{全降}
    \end{aligned}
\end{equation}

声调语言在世界上相当常见：汉藏语系、东南亚语言(泰语5个声调、越南语6个声调)、撒哈拉以南非洲的很多语言(约鲁巴语用声调区分词义)，都是声调语言。普通话四声可以用五度制([55]、[35]、[214]、[51])精确表示——这是\emph{调值}的离散编码。

\subsubsection{重音语言}

重音(响度/音高/时长的强调)在重音语言中起组织作用。重音可以是：
\begin{itemize}
    \item \textbf{固定重音}: 重音位置由规则固定(如波兰语几乎总是倒数第二音节)
    \item \textbf{自由重音}: 重音位置区分词义(如俄语 \cyr{м\'ука面粉 vs мук\'а痛苦})
    \item \textbf{音高重音}: 每个词只有一个高音调(日语、瑞典语)
\end{itemize}

英语名词重音规则(简化)：
\begin{equation}
    \text{重音落在倒数第三个音节，若倒数第二音节为重音节}
\end{equation}

\subsubsection{日语是声调还是重音语言?}

日语属于\emph{音高重音}(pitch accent)语言：每个词最多一个高调峰，单词之间通过调峰位置区分：\textit{はし}(chopsticks，头高型) vs \textit{はし}(bridge，尾高型)。这是介于声调语言与重音语言之间的类型，也是跨语言音系差异的极好案例。

\subsection{韵律与语调}

\begin{itemize}
    \item \textbf{音节/节拍等时性}: 英语是重音计时(stress-timed)，汉语、日语是音节/莫拉计时(syllable-timed / mora-timed)
    \item \textbf{语调(Intonation)}: 句子的整体音高曲线表达语气(陈述降调、疑问升调)。语调在所有语言中都存在，但功能有别：汉语中语调与字调叠加
\end{itemize}

\begin{equation}
    \text{普通话: 字调(词内) + 语调(句内) 叠加，语调只改变音高基线}
\end{equation}

\subsection{优选论(Optimality Theory)}

\subsubsection{OT框架}

优选论(Prince \& Smolensky, 1993)用\emph{约束排序 + 并行评估}取代串行规则：
\begin{definition}[OT框架]
OT由以下三部分组成：
\begin{itemize}
    \item GEN: 生成候选输出集合(所有可能的表层形式)
    \item CON: 普遍约束集合(标记性约束、忠实性约束)
    \item EVAL: 根据约束排序评估候选，选出违反最少(违反最高层级约束最少)的候选
\end{itemize}
\end{definition}

\begin{equation}
    \text{最优输出} = \arg\min_{c \in \text{GEN}} \sum_{i} w_i \cdot \text{Violations}(c, C_i)
\end{equation}

\subsubsection{约束的类型}

\begin{itemize}
    \item \textbf{标记性约束(Markedness)}: 禁制复杂/难发音的结构(如 \textsc{No-Coda} 禁止音节尾辅音、\textsc{Onset} 要求音节有起音)
    \item \textbf{忠实性约束(Faithfulness)}: 要求输出与输入一致(如 \textsc{Max} 禁止删音、\textsc{Dep} 禁止插音)
\end{itemize}

语言差异在OT中表现为\emph{约束排序的不同}：每种语言是普遍约束集合的一种排序。

\begin{example}
英语允许音节尾辅音(\emph{cat})，因为忠实性约束 \textsc{Max} 排序高于标记性约束 \textsc{No-Coda}；而部分语言(如某些波利尼西亚语)禁止音节尾辅音，是因为 \textsc{No-Coda} 排序更高。同一个普遍约束集合，不同的排序 = 不同的语言音系。这为"跨语言差异的形式化"提供了最优雅的范例：\emph{差异 = 参数(排序)选择}。
\end{example}

\subsubsection{OT与跨语言差异}

OT的意义在于把语言差异归约为\emph{排序}：假定所有语言共享同一个约束集合(普遍语法)，不同语言只是对约束的相对优先级排列不同。这与原则与参数理论(语法参数的二进制取值)异曲同工，共同体现了"普遍原则 + 参数化差异"的跨语言比较纲领。

\subsection{语音变化与历时音系学}

\subsubsection{常见的音变}

\begin{itemize}
    \item \textbf{同化(Assimilation)}: 音素变得与邻近音相似(英语复数清浊同化)
    \item \textbf{元音大转移}: 英语15-17世纪的元音系统性移位(现代英语拼写与发音不一致的原因)
    \item \textbf{音位化/去音位化}: 音位变体分裂为新音位，或音位对合并
\end{itemize}

\subsubsection{汉语语音史}

中古汉语到现代普通话发生了系统的辅音、元音、声调演变，音韵学用等韵图精确记录了这种演变。汉语的声调正是从早期无声调、经过声母清浊消失而"补偿"产生的。

\subsection{本章小结}

语音学测量声音，音系学组织声音。区别特征、音系规则、音位与音位变体、音节模板、声调与重音、优选论共同构成了研究"语音如何系统化"的完整工具箱。跨语言在音系上的差异极为显著且容易感知：音位库存、送气对立、音节结构、声调有无，都让不同语言"听起来"截然不同。而优选论等理论告诉我们：这些差异仍然来自共享原则的参数化排列。下一章进入计算语言学，看看语言知识如何变成可计算的系统。

\newpage
\section{计算语言学}

\subsection{概述}

计算语言学(Computational Linguistics, CL)使用计算方法研究语言，是自然语言处理(NLP)的理论基础。它既是语言学的分支(用计算模型验证语言理论)，也是计算机科学的分支(让计算机理解和生成语言)。

\begin{equation}
    \text{CL} = \text{语言学理论} \times \text{算法与统计模型}
\end{equation}

本书前几章讨论的语言学理论——形态学、句法、语义、音系——在这里都转化为可计算的模型：FST形态分析、PCFG句法分析、组合语义、HMM词性标注。

\subsection{语言模型}

\subsubsection{N-gram模型}

语言模型给任意字符串赋予概率，用于回答"这句话有多自然"。联合概率按链式法则展开：
\begin{equation}
    P(w_1, \ldots, w_n) = \prod_{i=1}^{n} P(w_i | w_1, \ldots, w_{i-1})
\end{equation}

但由于历史词序列太长，实际无法估计，因此引入\emph{马尔可夫假设}(n-gram近似)：
\begin{equation}
    P(w_i | w_1, \ldots, w_{i-1}) \approx P(w_i | w_{i-n+1}, \ldots, w_{i-1})
\end{equation}

即"当前词只依赖最近的 $n-1$ 个词"。bigram($n=2$)是最常用的一种。

\subsubsection{最大似然估计}

\begin{equation}
    P(w_i | w_{i-1}) = \frac{C(w_{i-1}, w_i)}{C(w_{i-1})}
\end{equation}
其中 $C(\cdot)$ 是语料中的计数。若某个bigram未出现，则概率为零——这会导致整个句子概率为零，因此需要平滑。

\subsubsection{平滑技术}

平滑技术为未见事件分配非零概率。Add-k平滑：
\begin{equation}
    P_{\text{add-k}}(w_i | w_{i-1}) = \frac{C(w_{i-1}, w_i) + k}{C(w_{i-1}) + k \cdot |V|}
\end{equation}

更先进的平滑方法包括：Good-Turing估计、Kneser-Ney平滑(目前最常用)以及回退(back-off)模型。平滑的本质是对"零计数"事件按频率衰减重新分配概率，同时保持总和为1。

\subsubsection{语言模型与语言差异}

语言模型直接反映语言的统计规律：训练自不同语言的语言模型，其n-gram分布差异巨大。例如汉语没有词边界(分词是预处理难题)、日语n-gram以"假名/形态素"为单位、形态丰富的语言(俄语、芬兰语)词形众多导致数据稀疏。这些差异催生了"子词(tokenization)"技术(如BPE)，成为现代多语言NLP的标准做法。

\subsection{词性标注}

\subsubsection{隐马尔可夫模型(HMM)}

词性标注的任务是：给定词序列，推断最可能的词性序列。HMM假设：
\begin{equation}
    P(t_1, \ldots, t_n | w_1, \ldots, w_n) = \prod_{i=1}^{n} P(t_i | t_{i-1}) \cdot P(w_i | t_i)
\end{equation}

\begin{itemize}
    \item $P(t_i | t_{i-1})$: 转移概率(词性之间的过渡)
    \item $P(w_i | t_i)$: 发射概率(某个词性下出现某个词)
\end{itemize}

HMM是一个生成模型：先按转移概率生成词性序列，再按发射概率生成词序列。标注就是求使 $P(t | w)$ 最大的词性序列，用贝叶斯法则：
\begin{equation}
    \hat{t} = \arg\max_t \prod_{i=1}^{n} P(t_i | t_{i-1}) \cdot P(w_i | t_i)
\end{equation}

\subsubsection{Viterbi算法}

HMM的解码用动态规划。定义 $v_t(j)$ 为"时刻 $t$ 处于状态 $j$ 且路径最优的最大概率"：
\begin{equation}
    v_t(j) = \max_{i} [v_{t-1}(i) \cdot a_{ij}] \cdot b_j(o_t)
\end{equation}
其中 $a_{ij}$ 是转移概率，$b_j(o_t)$ 是发射概率。同时用回溯指针记录最优路径。Viterbi时间复杂度 $O(T \cdot |S|^2)$，是NLP中最经典的算法之一。

\subsubsection{词性标注的跨语言差异}

不同语言的词性标注难度不同：
\begin{itemize}
    \item 英语: 形态简单，歧义主要来自词形相同(like可以是动词/介词/形容词)
    \item 汉语: 无词边界，分词与标注联合进行；词性相对稳定
    \item 俄语/德语: 屈折丰富，词形本身携带大量信息，但词形变化大、数据稀疏
    \item 日语: 词与词之间无空格，需先形态分析(分词)再标注
\end{itemize}

\subsection{词向量表示}

\subsubsection{分布语义假说}

"词的意义由其使用语境决定"(Firth, 1957)。分布语义把词表示为高维向量，使语义相似性转化为向量距离。

\subsubsection{Word2Vec}

Word2Vec(Mikolov et al., 2013)用神经网络从语料学习词向量，两种架构：
\begin{equation}
    \begin{aligned}
    &\text{Skip-gram: } \max \sum_{t} \sum_{-c \leq j \leq c, j \neq 0} \log P(w_{t+j} | w_t) \\
    &P(w_O | w_I) = \frac{\exp({v'_{w_O}}^\top v_{w_I})}{\sum_{w \in V} \exp({v'_w}^\top v_{w_I})}
    \end{aligned}
\end{equation}

Skip-gram用中心词预测上下文词，通过softmax输出概率。优化用负采样加速(不遍历整个词表)。

\subsubsection{词向量类比}

词向量编码了语义关系(词向量空间中的平移向量对应语义关系)：
\begin{equation}
    \text{king} - \text{man} + \text{woman} \approx \text{queen}
\end{equation}

\begin{equation}
    \text{king} - \text{man} \approx \text{queen} - \text{woman} \quad (\text{性别关系的平行位移})
\end{equation}

\subsubsection{多语言词向量}

多语言词向量把不同语言映射到同一向量空间，实现跨语言语义对齐：
\begin{equation}
    \text{英语} \xrightarrow{\text{线性变换}} \text{西班牙语向量空间}
\end{equation}
这种对齐(如MUSE、VecMap)使"语言之间的差异"可以被定量刻画：同一概念在不同语言中的向量相近，而语言特有的结构差异反映在对齐误差上。

\subsection{句法分析}

\subsubsection{依存分析}

依存分析预测每个词的中心词与关系：
\begin{equation}
    P(y|x) = \frac{\exp(w \cdot \Phi(x, y))}{\sum_{y'} \exp(w \cdot \Phi(x, y'))}
\end{equation}
其中 $\Phi(x,y)$ 是特征向量(词形、词性、距离等)，$w$ 是权重。这是最大熵/softmax分类器，可由带标注的树库训练。

\subsubsection{转移系统}

主流依存分析器(如Greedy transition-based parser)把句法分析建模为状态转移：
\begin{equation}
    \begin{aligned}
    &\text{SHIFT: } (\sigma, i|Q, A) \Rightarrow (\sigma|i, Q, A) \\
    &\text{LEFT-ARC: } (\sigma|i, j|Q, A) \Rightarrow (\sigma|i, Q, A \cup \{(j, r, i)\}) \\
    &\text{RIGHT-ARC: } (\sigma|i, j|Q, A) \Rightarrow (\sigma|i|j, Q, A \cup \{(i, r, j)\})
    \end{aligned}
\end{equation}

Arc-Eager系统维护栈 $\sigma$、队列 $Q$ 和已建立弧的集合 $A$，每次选择一个操作(移进/建左弧/建右弧)。贪心解析器用分类器选择下一个动作，复杂度 $O(n)$。

\subsubsection{跨语言依存分析: Universal Dependencies}

Universal Dependencies(UD)项目为100+语言标注统一的依存关系与词性体系。它的价值在于：\emph{用同一套标注体系描述世界语言}，使得跨语言句法研究(依存距离、语序倾向、树结构统计)成为可能。这正是"跨语言差异的形式化比较"的计算基础设施。

\subsection{Transformer与注意力}

\subsubsection{自注意力机制}

注意力机制让模型动态聚焦输入的相关部分。缩放点积注意力：
\begin{equation}
    \text{Attention}(Q, K, V) = \text{softmax}\left(\frac{QK^\top}{\sqrt{d_k}}\right)V
\end{equation}

\begin{equation}
    \begin{aligned}
    Q &= XW_Q \\
    K &= XW_K \\
    V &= XW_V
    \end{aligned}
\end{equation}

其中 $Q$(查询)、$K$(键)、$V$(值)是输入 $X$ 经过不同权重矩阵投影的结果，除以 $\sqrt{d_k}$ 防止softmax饱和。

\subsubsection{多头注意力}

多个注意力头并行捕捉不同关系：
\begin{equation}
    \text{MultiHead}(X) = \text{Concat}(\text{head}_1, \ldots, \text{head}_h)W_O
\end{equation}
\begin{equation}
    \text{head}_i = \text{Attention}(XW_Q^i, XW_K^i, XW_V^i)
\end{equation}

研究表明不同注意力头分别捕捉句法依存(主谓关系)、共指、长距离依赖等——注意力模式部分再现了语言学知识(依存树、语法角色)。Transformer也因此被称为"隐式学习语法的模型"。

\subsubsection{位置编码}

自注意力本身对位置不敏感，需要加入位置编码：
\begin{equation}
    PE_{(pos, 2i)} = \sin(pos / 10000^{2i/d})
\end{equation}
\begin{equation}
    PE_{(pos, 2i+1)} = \cos(pos / 10000^{2i/d})
\end{equation}

\subsection{大语言模型}

\subsubsection{语言建模目标}

大语言模型(LLM)的训练目标是自回归语言建模：最大化给定前文的条件下预测下一词的概率。
\begin{equation}
    \mathcal{L} = -\sum_{t=1}^{T} \log P(w_t | w_{<t}; \theta)
\end{equation}

这个"预测下一个词"的目标虽然简单，却迫使模型学习语法、语义、常识和世界知识。

\subsubsection{缩放定律}

模型性能随参数量、数据量、计算量遵循幂律(缩放定律, Scaling Laws)：
\begin{equation}
    L(N, D) = \frac{A}{N^\alpha} + \frac{B}{D^\beta} + E
\end{equation}
其中 $N$ 是参数量，$D$ 是训练数据量，$E$ 是数据不可约损失，$\alpha, \beta$ 是幂律指数。

\subsubsection{多语言模型}

现代LLM(如GPT系列)同时训练几十上百种语言，其多语言能力为"语言差异"提供了新的计算视角：
\begin{itemize}
    \item 模型内部形成了相对"语言中立"的语义表示(语言共享意义空间)
    \item 低资源语言受益于跨语言迁移(从数据多的语言学到模式)
    \item 模型对相近语言(英德、中韩)的迁移效果显著优于远缘语言
\end{itemize}
这从计算角度验证了语义学一章的结论：\emph{语义结构跨语言共享，表层编码各异}。

\subsection{机器翻译}

\subsubsection{统计机器翻译}

基于词的翻译模型：
\begin{equation}
    P(e | f) = \sum_{a} P(f | e, a) \cdot P(e)
\end{equation}
其中 $a$ 是词对齐。经典IBM模型用期望最大化(EM)估计对齐概率。

\subsubsection{神经机器翻译}

Transformer序列到序列模型直接学习 $P(e|f)$，端到端训练。注意力让解码器动态对齐源语言的对应片段。机器翻译的性能与语言对之间的距离相关：近缘语言对(英法)容易，远缘类型(英日、中英)难——语言差异的量化体现在翻译难度上。

\subsection{评估指标}

\subsubsection{困惑度(Perplexity)}

困惑度衡量语言模型的不确定度：
\begin{equation}
    \text{PPL}(W) = \sqrt[N]{\frac{1}{P(w_1, \ldots, w_N)}} = 2^{-\frac{1}{N}\sum_{i=1}^{N}\log_2 P(w_i | w_{<i})}
\end{equation}
困惑度越低越好(模型越确定)。注意PPL只能比较同一token化方案下的模型。

\subsubsection{BLEU分数}

BLEU衡量机器翻译与参考译文的 n-gram 重合度：
\begin{equation}
    \text{BLEU} = \text{BP} \cdot \exp\left(\sum_{n=1}^{N} w_n \log p_n\right)
\end{equation}
其中 $p_n$ 是 $n$-gram 精确率，BP是惩罚过短译文的"简短惩罚"。

\subsubsection{跨语言评估的挑战}

不同语言的评估难度不同：形态丰富的语言(俄语)BLEU对词形变化敏感；汉语、日语无空格分词影响token化；这些技术性差异本身也是"语言差异"的一部分。

\subsection{计算语言学与语言学的相互反哺}

计算语言学不只是应用语言学知识，它也反哺理论语言学：
\begin{itemize}
    \item \textbf{验证理论}: 依存分析器、词向量在数据上检验了语言学假说(主谓优先、距离倾向)
    \item \textbf{发现模式}: 大规模多语言模型揭示了语言共性(如依存距离的普遍最短化倾向)
    \item \textbf{量化差异}: 语料库方法把"语言差异"从定性描述变成可测度的数值(语素密度、语序倾向、依存距离)
\end{itemize}

\subsection{本章小结}

计算语言学把前几章的语言学理论转化为可计算的模型：语言模型刻画词汇的统计规律，HMM和Viterbi做词性标注，Word2Vec与Transformer学习语义表示，依存分析器与UD项目做跨语言句法分析，LLM通过下一个词预测隐式习得语法。计算视角下，"语言差异"成为可以定量测量的对象——而所有现代系统都验证了一个结论：语言共享深层结构，差异在表层参数与统计分布。

至此，我们已经系统学习了语言学各核心领域的理论。最后一章把这些视角整合起来，专门讨论\emph{不同语言之间的差异}：差异有哪些维度、如何形式化、以及为什么语言既"各不同"又"都一样"。

\newpage
\section{跨语言比较:不同语言的差异}

这是本书的总结章，也是核心关注的焦点。前面各章从不同层面建立了语言学的工具，本章把它们整合起来，系统地回答一个问题：\emph{不同的语言到底差在哪里?} 我们会看到，语言的差异是多维度、系统性的，但同时受到强烈约束——语言在\emph{共享深层结构}的前提下，以\emph{有限的方式参数化}地彼此不同。

\subsection{比较语言的维度}

\begin{table}[H]
\centering
\caption{语言差异的比较维度}
\begin{tabulary}{\textwidth}{@{}CCC@{}}
\toprule
\textbf{维度} & \textbf{内容} & \textbf{相关章节} \\
\midrule
语音/音系 & 音位库存、声调、音节结构 & 音系学与语音学 \\
词法 & 形态类型、格/数/性/时体系统 & 词法 \\
句法 & 语序、头参数、话题/主语结构 & 句法 \\
语义-词汇 & 词义切分、语法化、敬语 & 语义学 \\
文字 & 文字类型、正字法 & (本章) \\
语用 & 话题结构、信息结构、间接言语行为 & (本章) \\
\bottomrule
\end{tabulary}
\end{table}

比较方法上，除了定性描述，现代类型学用\emph{量化指标}比较语言：语素密度、基本语序分布、音位库存大小、依存距离等，都可以从语料库中自动计算。语言之间的"差异距离"也因此可以定义为一个多维向量空间中的距离。

\subsection{语音层面的差异}

\subsubsection{音位库存的规模差异}

不同语言的音位数量差异巨大。根据UPSID数据库：
\begin{itemize}
    \item \textbf{音位少的语言}: 夏威夷语(约13个音位)、皮拉罕语
    \item \textbf{音位中等的语言}: 英语(约44个音位)、汉语普通话(约21辅音+6元音+4声调)
    \item \textbf{音位多的语言}: !Xóõ语(非洲，辅音音位数百个，含复杂的搭嘴音)、高加索语言(辅音密集)
\end{itemize}

元音系统差异尤其直观：日语5个元音、英语15个左右、越南语11个、某些德语方言有20个以上。

\subsubsection{送气对立:一音两用 vs 两音对立}

\begin{itemize}
    \item 英语: 送气与不送气同属一个音位(变体分布，不区分意义)
    \item 印地语、泰语: 送气与不送气是不同音位(\emph{p}al vs \emph{p}hal 区分词义)
    \item 汉语普通话: /p/ /p\textsuperscript{h}/ 送气对立区分词义("爸" vs "怕")
\end{itemize}

\begin{equation}
    \text{英语: } [p] \sim [p^h] = /p/\ (\text{同一音位}) \qquad \text{汉语: } /p/ \neq /p^h/
\end{equation}

\subsubsection{声调的有无}

\begin{itemize}
    \item \textbf{声调语言}: 汉语(4调)、泰语(5调)、越南语(6调)、约鲁巴语(3调)
    \item \textbf{音高重音语言}: 日语、瑞典语(词内调峰位置区分词义)
    \item \textbf{重音语言}: 英语、俄语、德语、西班牙语(重音位置或固定或自由)
\end{itemize}

英语母语者听汉语"妈/麻/马/骂"觉得是同一音节的不同"语气"，而汉语母语者听英语重音差异如 "REcord"(名词) / "reCORD"(动词) 也难以掌握——两种语言的韵律系统结构不同。

\subsubsection{音节结构差异}

\begin{table}[H]
\centering
\caption{音节结构对比}
\begin{tabulary}{\textwidth}{@{}CCCC@{}}
\toprule
\textbf{语言} & \textbf{最大音节} & \textbf{允许辅音丛} & \textbf{例} \\
\midrule
汉语普通话 & (C)V(N) & 否 & 天 tiān \\
日语 & (C)(y)V(N/Q) & 否 & 東京 tōkyō \\
英语 & CCCVCCCC & 是 & strengths \\
波兰语 & CCCCVCCCC & 是(极复杂) & bezwzględy \\
\bottomrule
\end{tabulary}
\end{table}

音节模板决定了"外来词改造"的方式：英语单词进入日语被拆成开音节(computer → コンピューター)，进入汉语被套进(C)V(N)模板并加上声调。

\subsection{词法层面的差异}

\subsubsection{形态类型谱系}

语言在"一个词携带多少语法信息"上差异极大，构成一个连续谱：
\begin{equation}
    \text{孤立(汉语)} \xrightarrow{} \text{黏着(日语、土耳其语)} \xrightarrow{} \text{屈折(俄语、德语)} \xrightarrow{} \text{复综(因纽特语)}
\end{equation}

\begin{example}[同一句"我吃"在不同语言的形态差异]
\begin{align}
    \text{汉语: } & \text{我吃} \quad (\text{零形态，靠语序和代词}) \\
    \text{英语: } & \text{I eat} \quad (\text{少数屈折: 第三人称-s}) \\
    \text{日语: } & \text{食べる} \quad (\text{黏着词缀堆叠: 食べ-る}) \\
    \text{俄语: } & \text{\cyr{Я ем}} \quad (\text{动词变位融合人称-数-时}) \\
    \text{拉丁语: } & \text{ed\=o} \quad (\text{一个词尾融合所有语法信息})
\end{align}
\end{example}

\subsubsection{名词的性:语义 vs 语法}

\begin{itemize}
    \item \textbf{无语法性}: 汉语、日语、土耳其语、英语(基本没有)
    \item \textbf{两性}: 法语、西班牙语、意大利语(阴阳)
    \item \textbf{三性}: 德语、俄语、拉丁语(阴阳中)
    \item \textbf{多类}: 斯瓦希里语(名词类别系统，约10类)、班图语言
\end{itemize}

语法性(grammatical gender)与自然性别不一定对应：德语 \emph{Mädchen}(女孩)是中性的。跨语言比较中，同一个物体可以有不同的性：月亮在拉丁语、俄语是阴性，在德语是阳性。

\subsubsection{格的系统:从语序到词形}

\begin{table}[H]
\centering
\caption{格系统对比}
\begin{tabulary}{\textwidth}{@{}CCCC@{}}
\toprule
\textbf{语言} & \textbf{名词格} & \textbf{表达方式} & \textbf{例(宾格)} \\
\midrule
汉语 & 无 & 语序 & 猫吃鱼 / 鱼吃猫 \\
英语 & 仅代词保留 & 语序+残存代词格 & him \\
日语 & 有(助词) & 后置助词 & 魚を食べる \\
德语 & 4格 & 冠词+名词词尾 & den Fisch \\
俄语 & 6格 & 名词词尾 & \cyr{р\'ыбу} \\
芬兰语 & 15格 & 名词词尾 & kalaa \\
\bottomrule
\end{tabulary}
\end{table}

格与语序的关系是类型学的经典发现：\emph{格标记越丰富，语序越自由}。俄语、拉丁语格标记完备，语序可以相当自由(靠格判断谁吃谁)；汉语无语标记，语序成为唯一手段("猫吃鱼"不能换序为"鱼吃猫")。

\begin{equation}
    \text{形态格标记强度} \propto \text{语序自由度}
\end{equation}

\subsubsection{动词的时体系统}

\begin{itemize}
    \item \textbf{英语}: 时态(现在/过去)+ 体态(进行/完成)通过形态组合，将来时用助动词
    \item \textbf{汉语}: 无时态形态，用时间词和体助词"了/着/过"表达
    \item \textbf{俄语}: 完成/未完成体是词汇-语法对(\cyr{читать读/прочитать读完})，过去时用词形
    \item \textbf{日语}: 时态(-た过去)+ 体态(-ている进行/持续)
\end{itemize}

\begin{example}["我读了书"的跨语言编码]
\begin{align}
    \text{英语: } & \text{I read the book} \quad (\text{时态形态}) \\
    \text{汉语: } & \text{我读了书} \quad (\text{体助词 "了"}) \\
    \text{日语: } & \text{本を読んだ} \quad (\text{时态词尾 -だ}) \\
    \text{俄语: } & \text{\cyr{Я прочитал книгу}} \quad (\text{完成体前缀 \cyr{про-}})
\end{align}
\end{example}

\subsection{句法层面的差异}

\subsubsection{基本语序}

\begin{table}[H]
\centering
\caption{语序差异及代表语言}
\begin{tabulary}{\textwidth}{@{}CCCC@{}}
\toprule
\textbf{语序} & \textbf{代表语言} & \textbf{例句(我吃苹果)} \\
\midrule
SVO & 英语、汉语、法语、俄语 & I eat an apple \\
SOV & 日语、土耳其语、印地语 & 私はりんごを食べる \\
VSO & 阿拉伯语、威尔士语 & \cyr{(书面阿拉伯语)أكُلُ تفاحة} \\
\bottomrule
\end{tabulary}
\end{table}

\subsubsection{头参数:语序的系统性}

中心词的位置是"一把钥匙"，一次决定多种结构的语序。比较英语(中心词在前)与日语(中心词在后)：

\begin{table}[H]
\centering
\caption{头参数的连锁效应}
\begin{tabulary}{\textwidth}{@{}CCC@{}}
\toprule
\textbf{结构} & \textbf{英语(中心词前)} & \textbf{日语(中心词后)} \\
\midrule
动词+宾语 & eat an apple & りんごを 食べる \\
介词+名词 & in the house & 家の 中(名词+助词) \\
名词+属格 & the book of John & ジョンの 本 \\
名词+关系从句 & the book that I read & 私が読んだ 本 \\
形容词+名词 & red car & 赤い 車 \\
\bottomrule
\end{tabulary}
\end{table}

注意头参数不是铁律：汉语是SVO(中心词前)却用后置结构("桌子上")，属于"混合型"。但作为一个统计规律，头参数解释了英语与日语、土耳其语在几乎所有结构上的镜像语序。

\subsubsection{主语优先 vs 话题优先}

\begin{itemize}
    \item \textbf{主语优先(Subject-prominent)}: 句子以"主语-谓语"为骨架。英语、日语、印欧语系语言
    \item \textbf{话题优先(Topic-prominent)}: 句子以"话题-说明"为骨架。汉语、朝鲜语、日语
\end{itemize}

汉语的"话题"可以不是动词的施事：
\begin{equation}
    \text{这本书，我已经读完了} \quad (\text{话题 "这本书" 是受事})
\end{equation}

英语必须把它改为被动或补出受事："I have already read this book" / "This book has been read"。话题优先语言允许大量"无主语句"，且话题可以直接接评论，这在主语优先语言中不合语法。

\subsubsection{Pro-省略(主语省略)}

\begin{itemize}
    \item \textbf{可省略主语语言(pro-drop)}: 西班牙语、意大利语、日语、汉语
    \item \textbf{不可省略主语语言(non-pro-drop)}: 英语、法语、德语
\end{itemize}

\begin{example}
\begin{align}
    \text{意大利语: } & \text{Vado. (我)去。} \\
    \text{日语: } & \text{行く。(我/你/他…)去。} \\
    \text{英语: } & \text{*Go. (不合语法，必须说 I go / You go)}
\end{align}
\end{example}

pro-drop源于丰富的动词屈折(意大利语动词变位已标明主语人称)或话题结构(日语、汉语靠语境确定主语)。英语动词形态简化(几乎只有-s)，失去省略主语的条件。

\subsubsection{疑问句的形成}

\begin{table}[H]
\centering
\caption{是非疑问句的形成方式}
\begin{tabulary}{\textwidth}{@{}CCC@{}}
\toprule
\textbf{方式} & \textbf{代表语言} & \textbf{例} \\
\midrule
助动词移位 & 英语 & Are you ready? \\
句末语气词 & 汉语、日语 & 你准备好了吗? \\
动词前置 & 波兰语(部分) & Czy jesteś gotowy? \\
词尾形态 & 日语 & 行きますか \\
\bottomrule
\end{tabulary}
\end{table}

\subsubsection{Wh-疑问词的位置}

\begin{itemize}
    \item \textbf{疑问词移动}: 英语(What did you buy \_?)、德语、法语
    \item \textbf{疑问词原位(Wh-in-situ)}: 汉语(你买了什么?)、日语、韩语
\end{itemize}

虽然表层差异大，但深层一致：英语的 \_ 位置和汉语的"什么"位置在逻辑上完全相同(都是宾语)，区别只在"是否把疑问词提前"。这就是"普遍结构 + 参数"的典范。

\subsubsection{关系从句}

\begin{itemize}
    \item \textbf{后置关系从句}: 英语、法语、德语(the book that I read)
    \item \textbf{前置关系从句}: 日语、汉语、朝鲜语(私が読んだ本、"我读的书")
\end{itemize}

关系从句位置与头参数一致：中心词在后语言把关系从句放在名词前。这也产生著名的处理难点——前置关系从句的中心名词在从句之后才出现，读者需要"悬置"信息。

\subsection{代词、敬语与礼貌系统}

\subsubsection{第二人称的复杂性差异}

\begin{itemize}
    \item \textbf{英语}: 只有一个 you(不分单复、不分敬俗，600年前有 thou/you 区分)
    \item \textbf{法语/西班牙语}: tu/vous, tú/usted(亲昵/礼貌二分)
    \item \textbf{汉语}: 你/您(礼貌)，普通话无复数区分(口语"你们/您们")
    \item \textbf{日语}: 多个代词 + 敬语系统(见下)
\end{itemize}

\subsubsection{敬语系统:日语的极端例子}

日语敬语(敬語)分为三种，是跨语言差异中最复杂的礼貌系统之一：
\begin{itemize}
    \item \textbf{尊敬語(そんけいご)}: 抬高对方(いらっしゃる = いる/行く/来る的尊敬形)
    \item \textbf{謙譲語(けんじょうご)}: 贬低自己(参る = 行く/来る的自谦形)
    \item \textbf{丁寧語(ていねいご)}: 礼貌体(です/ます)
\end{itemize}

\begin{equation}
    \text{先生が来る} \xrightarrow{\text{尊敬}} \text{先生がいらっしゃる}
\end{equation}

朝鲜语也有复杂的敬语阶系统(等级由说话双方的社会关系决定)，英语则几乎完全依赖词汇选择(polite/please)。敬语系统反映的是语言背后的社会结构——语言差异不只是结构差异，也是文化差异。

\subsection{文字系统的差异}

\subsubsection{文字类型学}

\begin{table}[H]
\centering
\caption{文字类型}
\begin{tabulary}{\textwidth}{@{}CCC@{}}
\toprule
\textbf{文字类型} & \textbf{表示单位} & \textbf{代表} \\
\midrule
字母文字(Alphabet) & 音素 & 拉丁字母、西里尔字母 \\
元音附标文字(Abjad/Abugida) & 辅音(+元音附标) & 阿拉伯文、希伯来文、天城文 \\
音节文字(Syllabary) & 音节 & 日文假名、切罗基文 \\
语素文字(Logographic) & 词/语素 & 汉字 \\
\bottomrule
\end{tabulary}
\end{table}

\subsubsection{汉字与拼音文字的结构差异}

\begin{itemize}
    \item \textbf{汉字}: 语素-音节文字，一字对应一个音节+一个语素；形声字占多数(意符+声符，如"清"= 水+青)
    \item \textbf{假名}: 日语音节文字(五十音)，与汉字混用(汉字表实义词、假名表语法词缀)
    \item \textbf{谚文}: 韩语字母文字，把音素拼成方块字(\cyr{ㄱ+ㅏ+ㅁ=감})
    \item \textbf{拉丁字母}: 音素文字，拼写与发音的关系因语言而异(意大利语几乎一一对应，英语拼写与发音严重脱节)
\end{itemize}

\subsubsection{文字与语言结构的互动}

文字类型并非随机：孤立语(汉语)用一字一语的语素文字效率高；形态丰富的语言(俄语)用音素文字便于表达词形变化(同一词根的不同变格直接可见)。日语混合使用汉字(实词)与假名(词缀)来应付黏着形态，是文字与语言类型"适配"的典型案例。

\subsection{词汇与语法化的差异}

\subsubsection{词义切分的不同}

语言以不同的方式切分概念世界：
\begin{equation}
    \text{英语: } \text{hand} \quad \text{arm} \qquad \text{日语: } \text{手(て)} \quad \text{腕(うで)} \quad \text{同域}
\end{equation}

颜色、亲属、时间、空间等领域都是词义切分差异的经典例子(颜色词的普遍进化层级)。

\subsubsection{语法化的差异}

实词虚化为语法成分的过程(语法化)在不同语言中路径不同：
\begin{equation}
    \text{汉语: } \text{"把"(手抓)} \rightarrow \text{处置介词(把书放)} \qquad \text{英语: } \text{"have"} \rightarrow \text{完成助动词}
\end{equation}

汉语"了"从"完结"义动词语法化为体助词，英语 be going to 从"移动"语法化为将来时标记——同样的语法化倾向，不同语言的实现路径不同。

\subsection{形式化:原则与参数理论}

\subsubsection{语言差异 = 参数取值}

乔姆斯基的原则与参数理论(P\&P)为语言差异提供形式化框架：普遍语法提供原则(各语言共享)，语言之间的差异是少量\emph{参数}的取值不同。

\begin{equation}
    \text{语言 } L = \langle \text{普遍原则 } P, \text{参数向量 } \Theta(L) \rangle
\end{equation}

\begin{table}[H]
\centering
\caption{参数化差异(示例)}
\begin{tabulary}{\textwidth}{@{}CCCCC@{}}
\toprule
\textbf{参数} & \textbf{英语} & \textbf{汉语} & \textbf{日语} & \textbf{西班牙语} \\
\midrule
头参数 & 中心词前 & (混合) & 中心词后 & 中心词前 \\
主语省略 & 否 & 是(话题) & 是 & 是 \\
疑问词 & 移动 & 原位 & 原位 & 移动 \\
格标记 & 弱 & 无 & 助词 & 中 \\
声调 & 无 & 有(4调) & 音高重音 & 无 \\
\bottomrule
\end{tabulary}
\end{table}

这个框架的意义在于：儿童习得语言时，不需要学习"整套语法"，只需要根据输入\emph{设定参数值}。语言差异被压缩为少数二进制开关——这是对"语言为何既相似又不同"最简洁的解释。

\subsubsection{跨语言共性与蕴含层级}

语言差异不是任意的，普遍共性约束着差异。格林伯格共性和蕴含层级在差异中建立秩序：
\begin{equation}
    \text{若 } VSO \Rightarrow \text{前置词}; \qquad \text{格系统按层级} \quad \text{主格} \prec \text{宾格} \prec \text{与格} \prec \ldots
\end{equation}

统计上，语言有强烈的"默认设置"倾向：SOV和SVO占绝大多数、绝大多数语言动词紧跟主语或宾语、依存距离普遍保持最短。差异被强烈约束在少数"参数空间"中。

\subsection{语言相对性:差异是否影响思维?}

\subsubsection{萨丕尔-沃尔夫假说}

语言差异是否影响说话者的思维方式？萨丕尔-沃尔夫假说(Sapir-Whorf Hypothesis)有强(语言决定思维)与弱(语言影响思维)两个版本。

\begin{itemize}
    \item \textbf{强版(语言决定论)}: 已被广泛否定——人类普遍可以表达任意概念，任何语言都能表达新观念
    \item \textbf{弱版(语言相对论)}: 有实验支持——语言习惯影响注意与记忆的方式
\end{itemize}

\subsubsection{相对论的实验证据}

\begin{itemize}
    \item 颜色: 有基本颜色词的语言，说话者对颜色范畴的辨别更快
    \item 空间: 某些原住民语言用绝对方向(东西南北)而非相对方向(左右)描述空间，说话者的空间记忆方式随之不同
    \item 性别: 语法性影响说话者对无生命物的联想(德语桥是阴性、西班牙语桥是阳性)
\end{itemize}

\subsubsection{现代计算的视角}

大语言模型的多语言能力表明：不同语言可以表达完全相同的语义内容(可互相翻译)，支持"意义共享"立场。语言差异主要影响\emph{编码的效率与习惯}，而非表达能力本身。

\subsection{语言差异的量化比较}

现代计算类型学把差异转化为可测量指标，用距离度量语言：

\begin{equation}
    \text{语言距离} \; d(L_1, L_2) = \sqrt{\sum_{i} w_i \cdot (F_i(L_1) - F_i(L_2))^2}
\end{equation}

其中 $F_i$ 是类型学特征(语序、格标记、声调有无等)，$w_i$ 是权重。语系距离(词汇相似度)、语音距离(音位对应)、类型距离(结构特征)共同刻画两种语言的"差异度"。例如英语与德语距离小(同族同类型)，英语与日语距离大(不同族、类型迥异：SVO/孤屈折 vs SOV/黏着)。

\begin{table}[H]
\centering
\caption{英日对比:差异的维度总览}
\begin{tabulary}{\textwidth}{@{}CCC@{}}
\toprule
\textbf{维度} & \textbf{英语} & \textbf{日语} \\
\midrule
语系 & 印欧-日耳曼 & (独立/争议) \\
形态 & 弱屈折 & 黏着 \\
语序 & SVO & SOV \\
介词 & 前置词 & 后置助词 \\
格标记 & 弱(语序) & 强(助词) \\
声调 & 无 & 音高重音 \\
主语省略 & 否 & 是 \\
疑问词 & 移动 & 原位 \\
敬语 & 无系统 & 复杂系统 \\
文字 & 拉丁字母 & 汉字+假名 \\
\bottomrule
\end{tabulary}
\end{table}

\subsection{本章小结}

不同语言之间的差异是多维度的，但也是系统性的、受约束的：
\begin{itemize}
    \item \textbf{差异在表层}: 语音、词法、语序、文字——语言在最直观的层面差异巨大
    \item \textbf{共藏在深层}: 论元结构、组合语义、递归句法、普遍约束——所有语言共享
    \item \textbf{差异可参数化}: 原则与参数理论把差异压缩为参数向量，优选论把音系差异归约为约束排序
    \item \textbf{差异可量化}: 计算类型学让"语言距离"成为可测量的数量
\end{itemize}

因此，"语言的差异"与"语言的共性"是同一枚硬币的两面：差异是表层参数的选择，共性是深层结构的约束。理解了这一点，就掌握了语言学的精髓——从汉语的孤立语形态、英语的SVO语序、日语的黏着与敬语、俄语的格系统、阿拉伯语的VSO与声调，到所有语言共享的递归性与组合原则。

本书到此结束。我们从语言学的历史出发，历经形式语言、形态、句法、语义、音系与计算语言学，最终回到语言学的根本问题：语言是什么，以及语言如何既丰富多彩又和谐统一。


\end{document}
